% ============================================================================
% HORIZON PROTOCOL WHITEPAPER
% Professional LaTeX Technical Specification
% ============================================================================
% v1.0 - December 2025
% Author: Alex Metelli
% License: BSL-1.1 (converts to GPL-3.0 on 2028-12-19)

\documentclass[12pt,a4paper]{report}

% ============================================================================
% ENCODING AND FONTS
% ============================================================================

\usepackage[T1]{fontenc}
\usepackage{lmodern}                            % High-quality Latin Modern fonts

% ============================================================================
% CORE PACKAGES: Mathematics, Graphics, Tables
% ============================================================================

\usepackage{amsmath,amssymb,mathtools}          % Advanced mathematics
\usepackage{graphicx}                           % Graphics support
\usepackage{xcolor}                             % Color support

% ============================================================================
% TABLE PACKAGES: Professional multi-page support
% ============================================================================

\usepackage{booktabs}                           % Professional table formatting
\usepackage{longtable}                          % Tables spanning multiple pages
\usepackage{tabularx}                           % Responsive column widths
\usepackage{array}                              % Enhanced column formatting

% ============================================================================
% FLOAT CONTROL: Prevent tables drifting from context
% ============================================================================

\usepackage[section]{placeins}                  % Force floats before next section

% ============================================================================
% CODE LISTINGS: Clear spec/function blocks
% ============================================================================

\usepackage{listings}

\lstset{
  basicstyle=\ttfamily\footnotesize,            % Reduced font size for dense blocks
  frame=single,
  breaklines=true,
  columns=fullflexible,
  keepspaces=true
}

% ============================================================================
% ALGORITHM PSEUDOCODE: Professional function specifications
% ============================================================================

\usepackage[ruled,linesnumbered]{algorithm2e}

% Custom keywords for protocol specs
\SetKwInOut{Input}{Input}
\SetKwInOut{Output}{Output}
\SetKwInOut{Preconditions}{Preconditions}
\SetKwInOut{Postconditions}{Postconditions}
\SetKwInOut{StateChanges}{State Changes}
\SetKwInOut{Events}{Events}
\SetKwInOut{Reverts}{Reverts}
\SetKwFunction{FDeposit}{deposit}
\SetKwFunction{FRedeem}{redeem}
\SetKwFunction{FMintPY}{mint\_py}
\SetKwFunction{FRedeemPY}{redeem\_py}
\SetKwFunction{FClaimInterest}{claim\_interest}
\SetKwFunction{FSwap}{swap}
\SetKwFunction{FMint}{mint}
\SetKwFunction{FBurn}{burn}
\SetKw{Return}{return}
\SetKw{Revert}{revert}

% ============================================================================
% TYPOGRAPHY AND PAGINATION
% ============================================================================

\usepackage{microtype}                          % Advanced microtypography
\usepackage{needspace}                          % Prevent orphan headings
\usepackage{etoolbox}                           % Hook system
\usepackage{setspace}                           % Line spacing control
\usepackage{geometry}                           % Page layout

% Page geometry
\geometry{
  a4paper,
  left=1.1in,
  right=1.1in,
  top=1.1in,
  bottom=1.1in
}

% Line spacing (1.2 reduces eye fatigue on long documents)
\setstretch{1.2}

% Paragraph ergonomics: block paragraphs for specs (reduces fatigue)
\setlength{\parindent}{0pt}
\setlength{\parskip}{6pt plus 2pt minus 1pt}

% Prevent orphan headings (force minimum space before)
\pretocmd{\chapter}{\Needspace{8\baselineskip}}{}{}
\pretocmd{\section}{\Needspace{6\baselineskip}}{}{}
\pretocmd{\subsection}{\Needspace{5\baselineskip}}{}{}
\pretocmd{\subsubsection}{\Needspace{4\baselineskip}}{}{}

% ============================================================================
% SECTIONING TYPOGRAPHY: Explicit heading hierarchy
% ============================================================================

\usepackage{titlesec}

% Chapter format with subtle rule separator
\titleformat{\chapter}
  {\normalfont\LARGE\bfseries}
  {\thechapter}{1em}{}
  [\vspace{6pt}\titlerule]

\titlespacing*{\chapter}
  {0pt}{20pt}{30pt}

\titleformat{\section}
  {\normalfont\Large\bfseries}
  {\thesection}{1em}{}

\titlespacing*{\section}
  {0pt}{16pt}{14pt}

\titleformat{\subsection}
  {\normalfont\large\bfseries}
  {\thesubsection}{1em}{}

\titlespacing*{\subsection}
  {0pt}{12pt}{10pt}

\titleformat{\subsubsection}
  {\normalfont\normalsize\bfseries\itshape}
  {\thesubsubsection}{1em}{}

% ============================================================================
% TABLE OF CONTENTS IMPROVEMENTS
% ============================================================================

\usepackage{tocloft}
\setlength{\cftchapnumwidth}{2.5em}             % Prevent number collisions

% ============================================================================
% TABLE STYLING AND SPACING
% ============================================================================

\renewcommand{\arraystretch}{1.15}              % Table row spacing
\setlength{\tabcolsep}{6pt}                     % Column separation

% Equation spacing
\setlength{\abovedisplayskip}{10pt}
\setlength{\belowdisplayskip}{10pt}
\setlength{\abovedisplayshortskip}{8pt}
\setlength{\belowdisplayshortskip}{8pt}

% ============================================================================
% BUILD STABILITY: Prevent warnings and allow equation breaks
% ============================================================================

\sloppy                                         % Prevent overfull hbox warnings
\allowdisplaybreaks                             % Allow multi-line equations to break

% ============================================================================
% UNICODE CHARACTER SUPPORT
% ============================================================================

\usepackage{newunicodechar}

% Mathematical symbols
\newunicodechar{≥}{\ensuremath{\geq}}
\newunicodechar{≤}{\ensuremath{\leq}}
\newunicodechar{≠}{\ensuremath{\neq}}
\newunicodechar{→}{\ensuremath{\to}}
\newunicodechar{∞}{\ensuremath{\infty}}
\newunicodechar{×}{\ensuremath{\times}}
\newunicodechar{⋅}{\ensuremath{\cdot}}
\newunicodechar{∎}{\ensuremath{\blacksquare}}

% Greek letters
\newunicodechar{τ}{\ensuremath{\tau}}
\newunicodechar{φ}{\ensuremath{\varphi}}
\newunicodechar{π}{\ensuremath{\pi}}

% ============================================================================
% HYPERLINKS AND PDF METADATA
% ============================================================================

\usepackage[bookmarksnumbered=true]{hyperref}

\hypersetup{
  pdftitle={Horizon Protocol: A Yield Tokenization Protocol for Starknet},
  pdfauthor={Alex Metelli},
  pdfsubject={Technical Specification},
  pdfkeywords={yield tokenization, DeFi, Starknet, AMM, principal tokens, smart contracts, protocol specification},
  colorlinks=true,
  linkcolor=black,
  citecolor=black,
  urlcolor=blue,
  pdfborder={0 0 0}                             % Remove invisible link boxes
}

% PDF navigation: ensure subsections appear in outline
\setcounter{tocdepth}{3}
\setcounter{secnumdepth}{3}

% ============================================================================
% HEADERS AND FOOTERS
% ============================================================================

\usepackage{fancyhdr}
\pagestyle{fancy}
\fancyhf{}
\lhead{\small\nouppercase{\leftmark}}           % Chapter context (no forced uppercase)
\rhead{\small \thepage}
\renewcommand{\headrulewidth}{0.4pt}

% Increase headheight to prevent fancyhdr warnings
\setlength{\headheight}{14pt}
\addtolength{\topmargin}{-2pt}

% Plain style for chapter pages
\fancypagestyle{plain}{%
  \fancyhf{}
  \cfoot{\thepage}
  \renewcommand{\headrulewidth}{0pt}
}

% ============================================================================
% THEOREM ENVIRONMENTS (numbered by chapter for easier cross-referencing)
% ============================================================================

\usepackage{amsthm}

\theoremstyle{plain}
\newtheorem{theorem}{Theorem}[chapter]
\newtheorem{lemma}[theorem]{Lemma}
\newtheorem{corollary}[theorem]{Corollary}
\newtheorem{proposition}[theorem]{Proposition}

\theoremstyle{definition}
\newtheorem{definition}[theorem]{Definition}
\newtheorem{example}[theorem]{Example}

\theoremstyle{remark}
\newtheorem*{remark}{Remark}
\newtheorem*{notation}{Notation}

% ============================================================================
% DOCUMENT BEGIN
% ============================================================================

\begin{document}

% ============================================================================
% TITLE PAGE
% ============================================================================

\thispagestyle{empty}

\vspace*{50pt}

\begin{center}
  {\LARGE\bfseries Horizon Protocol}

  \vspace{20pt}

  {\Large A Yield Tokenization Protocol for Starknet}

  \vspace{40pt}

  {\large Technical Specification | Version 1.0}

  \vspace{60pt}

  {\normalsize December 2025}

  \vspace{30pt}

  {\normalsize Alex Metelli}

  \vspace{10pt}

  {\small \url{https://github.com/ametel01/horizon-starknet}}

  \vspace{40pt}

  {\footnotesize License: BSL-1.1 (converts to GPL-3.0 on 2028-12-19)}
\end{center}

\vfill

\clearpage

% ============================================================================
% ABSTRACT AND KEYWORDS
% ============================================================================

\chapter*{Abstract}
\addcontentsline{toc}{chapter}{Abstract}

\itshape
Horizon Protocol is a yield tokenization system deployed on Starknet that decomposes yield-bearing assets into separate Principal Tokens (PT) and Yield Tokens (YT). This separation enables three distinct financial strategies: fixed-yield positions through discounted PT acquisition, leveraged yield exposure via YT trading, and liquidity provision in the protocol's native PT/SY automated market maker.

The protocol operates on a fundamental conservation invariant: one unit of Standardized Yield token (SY) can be split into one PT and one YT before maturity, and this process is reversible. At maturity, each PT is redeemable for exactly one unit of the underlying SY, providing a guaranteed redemption value. YT holders receive all yield generated by the underlying asset until expiry, after which the token expires worthless.

Horizon employs a logit-based AMM curve specifically designed for time-decaying assets, ensuring PT price converges to 1 as maturity approaches. This specification provides formal definitions, mathematical proofs, economic analysis, and comprehensive risk assessment suitable for auditors, researchers, developers, and investors.
\normalfont

\vspace{12pt}

\noindent \textbf{Keywords:} yield tokenization, principal tokens, yield tokens, Starknet, DeFi, automated market makers, fixed-point arithmetic, smart contracts, protocol specification, mathematical proofs

\clearpage

% ============================================================================
% TABLE OF CONTENTS
% ============================================================================

\tableofcontents

\clearpage

% ============================================================================
% MAIN CONTENT
% ============================================================================

\chapter{Motivation \& Problem Statement}

% Lead paragraph establishing argument structure
This chapter establishes why yield tokenization is necessary, why existing solutions are inadequate, how tokenization addresses these gaps, and what Starknet-specific considerations shape the protocol's design.

\section{Yield Variability and Rate Uncertainty}

Yield-bearing tokens have become fundamental primitives in decentralized finance. Liquid staking derivatives (nstSTRK, sSTRK), wrapped staking tokens (wstETH), and lending protocol receipt tokens generate yield through staking rewards, lending interest, or protocol revenue distribution. However, these yields share a common characteristic: \textbf{variability}.

In practice, the same position can swing from low single-digit APY to double-digit APY over short horizons, driven by staking participation, protocol incentives, and utilization. For treasury management, this converts yield into an uncertain cash-flow stream rather than an allocatable budget line.

Consider a user holding nstSTRK (Nostra's staked STRK derivative). The yield depends on:
\begin{itemize}
    \item Starknet's staking reward rate, which fluctuates with network participation
    \item Nostra's protocol fee structure, subject to governance changes
    \item Market conditions affecting utilization and demand
\end{itemize}

This variability creates three requirements that existing primitives do not satisfy:

\begin{enumerate}
    \item \textbf{Forecastability}: Treasuries, DAOs, and institutional holders require predictable cash flows for budgeting and liability matching.

    \item \textbf{Rate-locking without collateral inefficiency}: Capital providers need fixed-rate exposure without overcollateralized debt positions or idle margin.

    \item \textbf{A market for yield views and hedging}: Users with convictions about future yield direction require an efficient mechanism to express these views or hedge existing exposure.
\end{enumerate}

\section{Why Existing Markets Do Not Solve Rate Certainty}

\subsection{Traditional Lending Protocols}

Protocols like Aave and Compound offer variable-rate lending and borrowing. While these markets provide yield, the rates fluctuate continuously based on utilization. A depositor earning 5\% today may earn 2\% tomorrow. These protocols solve capital efficiency but not rate certainty.

\subsection{Fixed-Rate Lending Protocols}

Protocols such as Notional Finance and Yield Protocol introduced fixed-rate borrowing and lending. However, these approaches face structural limitations:

\begin{center}
\begin{tabularx}{\textwidth}{>{\raggedright\arraybackslash}p{4cm} >{\raggedright\arraybackslash}X}
\toprule
\textbf{Limitation} & \textbf{Description} \\
\midrule
Governance-driven term listing & Fixed maturities require protocol governance to launch new terms \\
Liquidity fragmentation & Capital splits across multiple maturity pools \\
Overcollateralization & Borrowers lock collateral; capital sits idle \\
Limited secondary market & Users lock rates but cannot trade rate exposure \\
\bottomrule
\end{tabularx}
\end{center}

\subsection{Yield Tokenization as a Decomposition Primitive}

Yield tokenization decomposes a yield-bearing position into two separable claims: a principal claim (PT) on redemption value at expiry, and a yield claim (YT) on all pre-expiry appreciation. This decomposition addresses the requirements introduced above without introducing borrower-side collateral constraints.

\begin{enumerate}
    \item \textbf{Separation of Rights}: PT and YT can be held, traded, or used as collateral independently, enabling hedging and targeted exposure.

    \item \textbf{Continuous Price Discovery}: Market prices for PT and YT encode collective expectations about future yield over the remaining time to expiry.

    \item \textbf{Capital Efficiency}: Exposure is represented directly by transferable tokens rather than overcollateralized debt positions.

    \item \textbf{Composability}: PT and YT implement the ERC-20 interface pattern on Starknet and integrate with existing DeFi tooling.

    \item \textbf{Deterministic Maturity Handling}: PT mechanically converges to redemption value at expiry; settlement does not require term re-listing.
\end{enumerate}

\section{Starknet Constraints and Design Opportunities}

Implementing yield tokenization on Starknet requires careful numeric design (felt-based arithmetic, u256 emulation) but also benefits from predictable execution and native integration with Starknet yield assets.

\subsection{Constraints}

Cairo, Starknet's native smart contract language, operates on a field element (felt252) primitive with different arithmetic properties than EVM's uint256. The protocol requires:

\begin{itemize}
    \item Custom u256 handling for token amounts
    \item High-precision fixed-point math for rate calculations
    \item Explicit signed value handling for price differentials
\end{itemize}

The cubit 64.64 fixed-point library addresses these requirements, providing approximately 19 decimal digits of precision. All rate and price computations are defined in terms of deterministic fixed-point routines to avoid platform-dependent rounding.

\subsection{Integration with Starknet Yield Sources}

The protocol targets integration with Starknet's native yield-bearing tokens:

\begin{center}
\begin{tabular}{lll}
\toprule
\textbf{Token} & \textbf{Type} & \textbf{Yield Mechanism} \\
\midrule
nstSTRK & Liquid staking & Staked STRK + lending yield \\
sSTRK & Liquid staking & Staking rewards (derivative) \\
wstETH & Bridged staking & Ethereum staking rewards \\
\bottomrule
\end{tabular}
\end{center}

Each integration requires an exchange rate oracle to track the underlying token's value appreciation.

\subsection{Oracle Integration}

Pragma is a major oracle provider on Starknet. For bridged assets, oracle feeds avoid per-operation cross-chain dependencies. The protocol uses Pragma's TWAP (time-weighted average price) feeds to mitigate single-block manipulation.

\section{Assumptions and Threats to Guarantees}

The protocol's guarantees depend on the following external assumptions. These are dependencies the protocol cannot enforce on-chain.

\begin{center}
\begin{tabularx}{\textwidth}{>{\raggedright\arraybackslash}X >{\raggedright\arraybackslash}X}
\toprule
\textbf{Assumption} & \textbf{Failure Mode} \\
\midrule
Underlying asset valuation is monotonic & Non-monotonic exchange rate breaks YT accounting \\
Oracle liveness and bounded error & Stale feeds cause mispricing or incorrect settlement \\
Sufficient arbitrage participation & Invariants not economically enforced during stress \\
Consistent timestamp semantics & Edge cases around expiry boundary \\
Underlying protocol solvency & SY value impairment \\
\bottomrule
\end{tabularx}
\end{center}

These are external dependencies; the protocol cannot enforce them on-chain. Users must evaluate these dependencies when interacting with the protocol.


\chapter{System Overview}

\section{Protocol Architecture}

The Horizon Protocol is organized into layered components with clearly defined responsibilities and upgrade boundaries. This separation isolates external dependencies from core accounting logic and limits the surface area affected by upgrades.

\subsection{Layer Model}

\begin{center}
\begin{tabular}{lll}
\toprule
\textbf{Layer} & \textbf{Components} & \textbf{Upgrade Scope} \\
\midrule
External & Underlying yield tokens, Pragma Oracle & Outside protocol control \\
Core & SY, PT, YT, Market & Upgradeable via governed proxy \\
Periphery & Router, Factory, MarketFactory & Upgradeable via governed proxy \\
\bottomrule
\end{tabular}
\end{center}

\section{Token Flow}

The protocol implements a reversible asset flow centered around a deterministic split of standardized yield into principal and yield claims.

\subsection{Forward Flow (Deposit and Split)}

\begin{enumerate}
    \item A user deposits an underlying yield-bearing token into the SY contract.
    \item The SY contract mints a standardized yield token representing the deposited asset.
    \item The user deposits SY into the YT contract, which atomically mints paired PT and YT tokens such that $1\,\text{SY} = 1\,\text{PT} + 1\,\text{YT}$.
\end{enumerate}

\subsection{Reverse Flow (Redemption)}

\begin{enumerate}
    \item Before expiry, a user may redeem exactly one PT and one YT to recover one SY, reversing the split.
    \item At expiry, PT redeems deterministically 1:1 for SY, while YT expires with zero value.
    \item The user withdraws the underlying asset from the SY contract, receiving principal plus any accrued yield.
\end{enumerate}

\section{Token System}

The protocol defines four token types representing distinct economic claims. All tokens implement the ERC-20 interface pattern and impose no transfer restrictions.

\begin{center}
\begin{tabularx}{\textwidth}{l >{\raggedright\arraybackslash}X l l}
\toprule
\textbf{Token} & \textbf{Economic Role} & \textbf{Minting Mechanism} & \textbf{Transferability} \\
\midrule
SY & Standardized yield wrapper & Permissionless via deposit & Freely transferable \\
PT & Principal claim (redeemable at expiry) & Minted via YT contract & Freely transferable \\
YT & Yield claim (pre-expiry) & Minted via YT contract & Freely transferable \\
LP & Liquidity provider share & Minted by Market contract & Freely transferable \\
\bottomrule
\end{tabularx}
\end{center}

\section{Core Protocol Invariants}

The following invariants are safety-critical and must hold across all valid state transitions of the protocol.

\begin{enumerate}
    \item \textbf{Principal Protection}: Each PT redeems deterministically for exactly one unit of SY at maturity.
    \item \textbf{Conservation}: Prior to expiry, the combined value of PT and YT corresponds to one unit of SY.
    \item \textbf{Price Convergence}: The AMM enforces monotonic convergence of PT price toward redemption value as expiry approaches.
    \item \textbf{Yield Isolation}: All yield generated by the underlying asset accrues exclusively to YT holders prior to expiry.
\end{enumerate}


\chapter{Formal Model \& Notation}

This section formalizes the protocol's state, parameters, and mathematical invariants. All subsequent mechanism definitions rely on the notation and constants introduced here. We adopt a fixed-point arithmetic model (WAD scaling, $10^{18}$) to ensure precision across all computations.

\section{State Variables}

The protocol state is partitioned into market state, time state, and index state. Each variable has a well-defined domain that is enforced by construction.

\subsection{Market State}

\begin{center}
\begin{tabular}{llll}
\toprule
\textbf{Symbol} & \textbf{Type} & \textbf{Domain} & \textbf{Description} \\
\midrule
$S_{SY}$ & u256 & $[0, 2^{256})$ & SY reserve in market \\
$S_{PT}$ & u256 & $[0, 2^{256})$ & PT reserve in market \\
$L$ & u256 & $[0, 2^{256})$ & Total LP token supply \\
$r$ & u256 & $[0, 4.6 \cdot 10^{18}]$ & ln(implied rate) in WAD \\
$r_{root}$ & u256 & $[0, 10^{20}]$ & Scalar root parameter \\
$r_{anchor}$ & u256 & $[0, 10^{19}]$ & Rate anchor parameter \\
$\phi$ & u256 & $[0, 0.05 \cdot 10^{18}]$ & Fee rate in WAD \\
\bottomrule
\end{tabular}
\end{center}

\textbf{Assumption}: $S_{PT} + S_{SY} > 0$ at all times due to enforced minimum liquidity lock.

\subsection{Time State}

\begin{center}
\begin{tabular}{llll}
\toprule
\textbf{Symbol} & \textbf{Type} & \textbf{Domain} & \textbf{Description} \\
\midrule
$T$ & u64 & Unix timestamp & Expiry timestamp \\
$t$ & u64 & Unix timestamp & Current block timestamp \\
$\tau$ & u64 & $[0, T]$ & Time to expiry (seconds) \\
$\tau_y$ & u256 & $[0, \infty)$ & Time to expiry (years, WAD-scaled) \\
\bottomrule
\end{tabular}
\end{center}

\subsection{Index State}

\begin{center}
\begin{tabular}{llll}
\toprule
\textbf{Symbol} & \textbf{Type} & \textbf{Domain} & \textbf{Description} \\
\midrule
$I_{global}$ & u256 & $[10^{18}, \infty)$ & Global PY index (watermark) \\
$I_{user}$ & u256 & $[10^{18}, \infty)$ & User's last recorded index \\
\bottomrule
\end{tabular}
\end{center}

\section{Constants}

The protocol relies on the following global constants, which are fixed for the lifetime of a given market instance.

\textbf{Scaling Constants:}
\begin{itemize}
\item \textbf{WAD} ($W = 10^{18}$): Fixed-point scaling factor for all amounts and rates
\item \textbf{Seconds per year} ($Y = 31{,}536{,}000$): Time normalization constant for annualized rates
\end{itemize}

\textbf{Safety Bounds:}
\begin{itemize}
\item \textbf{Minimum liquidity} ($L_{min} = 1000$): Locked LP to prevent division by zero
\item \textbf{Max ln implied rate} ($r_{max} = 4.6 \cdot 10^{18}$): Caps implied rate at approximately 10,000\% APY
\item \textbf{Max proportion} ($p_{max} = 0.999 \cdot 10^{18}$): Prevents extreme pool imbalance
\item \textbf{Min proportion} ($p_{min} = 0.001 \cdot 10^{18}$): Prevents extreme pool imbalance
\end{itemize}

\section{Fixed-Point Arithmetic}

The protocol uses a dual-layer precision model to balance gas efficiency with computational accuracy.

\subsection{WAD Arithmetic (Interface Layer)}

All token amounts and external-facing values use WAD scaling:

\[
\operatorname{wad\_mul}(a, b) = \left\lfloor \frac{a \times b + W/2}{W} \right\rfloor
\]

\[
\operatorname{wad\_div}(a, b) = \left\lfloor \frac{a \times W + b/2}{b} \right\rfloor
\]

where $W = 10^{18}$ and rounding is to nearest integer.

All intermediate multiplications are performed using widened precision to prevent overflow. Rounding is symmetric and biased toward the nearest integer to minimize cumulative drift.

\subsection{cubit 64.64 (Internal Layer)}

Transcendental functions use 64.64 binary fixed-point representation with 64 integer bits and 64 fractional bits, providing approximately 19 decimal digits of precision.

The 64.64 format represents a number $x$ as:
\[
x_{\text{fixed}} = x \times 2^{64}
\]

\paragraph{Conversion Functions.}

\[
\operatorname{wad\_to\_fp}(x) = \left\lfloor \frac{x \times 2^{64}}{10^{18}} \right\rfloor
\]

\[
\operatorname{fp\_to\_wad}(x) = \left\lfloor \frac{x \times 10^{18}}{2^{64}} \right\rfloor
\]

\paragraph{High-Precision Functions.}

The library provides the following transcendental functions:

\begin{itemize}
\item $\operatorname{exp\_wad}(x)$: Exponential function, takes WAD input and returns WAD output
\item $\operatorname{exp\_neg\_wad}(x)$: Negative exponential for decay functions
\item $\operatorname{ln\_wad}(x)$: Natural logarithm, returns (value, is\_negative) tuple
\item $\operatorname{pow\_wad}(x, y)$: Power function computing $x^y$
\item $\operatorname{sqrt\_wad}(x)$: Square root function, WAD-normalized output
\end{itemize}

All transcendental functions clamp outputs on domain violation and return bounded values to preserve protocol safety.

\section{Time Model}

Time is measured in Unix timestamps (seconds since epoch):

\[
\tau = T - t
\]

where $T$ is expiry timestamp and $t$ is current block timestamp.

Time expressed in years (WAD-scaled):

\[
\tau_y = \frac{\tau \times W}{Y}
\]

For $t \geq T$, the protocol treats $\tau = 0$ and transitions to post-expiry redemption semantics.

\section{Core Equations}

This section defines the mathematical relationships governing AMM pricing, yield computation, and fee adjustment.

\subsection{Pool Proportion and Logit Transformation}

The proportion of PT in the pool determines market pricing:

\[
p = \frac{S_{PT}}{S_{PT} + S_{SY}}
\]

\textbf{Domain constraint}: $p_{min} \leq p \leq p_{max}$. Trades that would violate this constraint are rejected.

To map bounded pool proportions to an unbounded rate domain, the protocol uses a logit transformation:

\[
\operatorname{logit}(p) = \ln\left(\frac{p}{1-p}\right)
\]

Properties:
\begin{itemize}
    \item $\operatorname{logit}(0.5) = 0$
    \item $\operatorname{logit}(p) \to -\infty$ as $p \to 0$
    \item $\operatorname{logit}(p) \to +\infty$ as $p \to 1$
\end{itemize}

\subsection{Pricing Model}

Time-to-expiry is incorporated through a scalar that dampens price sensitivity as maturity approaches.

\paragraph{Rate Scalar.}

\[
r_{scalar} = \frac{r_{root} \times Y}{\tau}
\]

As $\tau \to 0$: $r_{scalar} \to \infty$, making prices insensitive to pool balance changes.

\paragraph{Exchange Rate.}

The PT/SY exchange rate derived from market state:

\[
\operatorname{exchangeRate} = \frac{\operatorname{logit}(p)}{r_{scalar}} + r_{anchor}
\]

\paragraph{Implied Rate.}

The annualized implied yield from the exchange rate:

\[
r = \ln(\operatorname{exchangeRate}) \times \frac{Y}{\tau}
\]

\textbf{Constraint}: $0 \leq r \leq r_{max}$. Values outside this range are clamped.

\paragraph{PT Price.}

PT price in terms of SY, derived from implied rate:

\[
P_{PT} = e^{-r \times \tau_y / W}
\]

where $\tau_y$ is time to expiry in WAD-scaled years.

Properties:
\begin{itemize}
    \item $P_{PT} < 1$ before expiry (trades at discount)
    \item $P_{PT} \to 1$ as $\tau \to 0$ (converges at maturity)
\end{itemize}

\paragraph{YT Price.}

YT price is derived directly from the conservation invariant:

\[
P_{YT} = 1 - P_{PT}
\]

Properties:
\begin{itemize}
    \item $P_{YT} > 0$ before expiry
    \item $P_{YT} \to 0$ as $\tau \to 0$ (expires worthless)
\end{itemize}

\subsection{Rate Anchor Dynamics}

The rate anchor maintains implied rate continuity across trades, preventing discontinuous price jumps.

\paragraph{Anchor Calculation.}

Given the last stored implied rate $r_{last}$ and current market state:

\[
\operatorname{exchangeRate}_{target} = e^{r_{last} \times \tau / Y}
\]

\[
r_{anchor} = \operatorname{exchangeRate}_{target} - \frac{\operatorname{logit}(p)}{r_{scalar}}
\]

\paragraph{Anchor Update.}

After each trade, the anchor is recalculated from the new implied rate to ensure smooth price transitions.

\subsection{Fee Model}

Swap fees decay exponentially toward expiry, reducing friction as price convergence becomes more certain.

\[
\phi_{adjusted} = e^{\ln(\phi_{root}) \times \tau / Y}
\]

where $\phi_{root}$ is the base fee rate (e.g., 0.001 WAD = 0.1\%).

\textbf{Requirement}: $0 < \phi_{root} < 1$ to ensure monotonic decay of swap fees.

\begin{center}
\begin{tabular}{ll}
\toprule
\textbf{Time to Expiry} & \textbf{Adjusted Fee (if base = 0.1\%)} \\
\midrule
1 year & 0.100\% \\
6 months & $\approx 0.070\%$ \\
1 month & $\approx 0.022\%$ \\
1 week & $\approx 0.006\%$ \\
At expiry & $\approx 0\%$ \\
\bottomrule
\end{tabular}
\end{center}

\subsection{Yield Index Model}

The PY index accounts for cumulative yield attributable to YT holders using a monotonic watermark mechanism.

\paragraph{Index Update.}

\[
I_{global} = \max(I_{global}, I_{current})
\]

where $I_{current}$ is the current exchange rate from the oracle.

The watermark ensures $I_{global}$ is monotonically non-decreasing. This prevents yield accounting errors during temporary rate decreases.

\paragraph{Interest Calculation.}

For a user with YT balance $B_{YT}$ and last recorded index $I_{user}$:

\[
\text{Interest} = B_{YT} \times \frac{I_{global} - I_{user}}{I_{user}}
\]

After claiming, the user's index is updated: $I_{user} \leftarrow I_{global}$

\section{Invariants}

The following invariants are safety-critical. Any violation constitutes a protocol failure and must be prevented by construction or rejected at runtime.

\begin{center}
\begin{tabularx}{\textwidth}{l >{\raggedright\arraybackslash}X l}
\toprule
\textbf{Invariant} & \textbf{Expression} & \textbf{Enforcement} \\
\midrule
Supply equality & $\text{PT.supply} = \text{YT.supply}$ & Contract logic \\
Conservation & PT + YT = SY (value) & Arbitrage \\
Index watermark & $I_{global}^{t+1} \geq I_{global}^{t}$ & Contract logic \\
Proportion bounds & $p_{min} \leq p \leq p_{max}$ & Trade rejection \\
Rate bounds & $0 \leq r \leq r_{max}$ & Clamping \\
Reserve positivity & $S_{SY} > 0 \land S_{PT} > 0$ & LP minimum lock \\
\bottomrule
\end{tabularx}
\end{center}


\chapter{Core Mechanism Design}

This section specifies the protocol's core operations with precise state transitions, revert conditions, and invariant preservation. Each mechanism is defined by its inputs, preconditions, state changes, events, returns, and revert conditions.

\section{SY Token Mechanism}

The Standardized Yield (SY) token wraps yield-bearing assets with a uniform interface. SY is a fixed 1:1 wrapper: one unit of underlying always mints one unit of SY, and one unit of SY always redeems one unit of underlying. The exchange rate tracked by the protocol is not the SY appreciation rate, but the underlying asset's index used for YT yield accounting.

\paragraph{Key distinction.} SY itself does not appreciate. The \texttt{exchange\_rate} (or \texttt{index}) tracks the underlying yield-bearing token's value growth and is used exclusively for calculating YT interest. This separation ensures SY remains a simple wrapper while yield accrual is handled by the YT mechanism.

\subsection{Deposit}

\begin{algorithm}[H]
\caption{SY Deposit}
\underline{function \textsc{Deposit}}$(receiver, amount) \rightarrow shares$\;
\Input{$receiver$: address, $amount$: u256}
\Output{$shares$: u256 (amount of SY minted)}
\Preconditions{
  $receiver \neq 0$\;
  $amount > 0$\;
  $\texttt{underlying.allowance}(caller, SY) \geq amount$\;
}
\BlankLine
$\texttt{underlying.balanceOf}(caller) \gets \texttt{underlying.balanceOf}(caller) - amount$\;
$\texttt{underlying.balanceOf}(SY) \gets \texttt{underlying.balanceOf}(SY) + amount$\;
$\texttt{SY.balanceOf}(receiver) \gets \texttt{SY.balanceOf}(receiver) + amount$\;
$\texttt{SY.totalSupply} \gets \texttt{SY.totalSupply} + amount$\;
\BlankLine
\textbf{emit} $\textsc{Deposit}(caller, receiver, amount, shares)$\;
\Return $shares \gets amount$\;
\BlankLine
\Reverts{
  $amount = 0$: ``zero amount''\;
  $receiver = 0$: ``zero address''\;
  allowance insufficient: ``ERC20: insufficient allowance''\;
  balance insufficient: ``ERC20: transfer exceeds balance''\;
}
\end{algorithm}

\subsection{Redeem}

\begin{algorithm}[H]
\caption{SY Redeem}
\underline{function \textsc{Redeem}}$(receiver, amount) \rightarrow assets$\;
\Input{$receiver$: address, $amount$: u256}
\Output{$assets$: u256 (amount of underlying returned)}
\Preconditions{
  $receiver \neq 0$\;
  $amount > 0$\;
  $\texttt{SY.balanceOf}(caller) \geq amount$\;
}
\BlankLine
$\texttt{SY.balanceOf}(caller) \gets \texttt{SY.balanceOf}(caller) - amount$\;
$\texttt{SY.totalSupply} \gets \texttt{SY.totalSupply} - amount$\;
$\texttt{underlying.balanceOf}(SY) \gets \texttt{underlying.balanceOf}(SY) - amount$\;
$\texttt{underlying.balanceOf}(receiver) \gets \texttt{underlying.balanceOf}(receiver) + amount$\;
\BlankLine
\textbf{emit} $\textsc{Redeem}(caller, receiver, amount, assets)$\;
\Return $assets \gets amount$\;
\BlankLine
\Reverts{
  $amount = 0$: ``zero amount''\;
  $receiver = 0$: ``zero address''\;
  balance insufficient: ``ERC20: burn exceeds balance''\;
}
\end{algorithm}

\subsection{Exchange Rate (Index)}

The exchange rate tracks the underlying asset's value appreciation and is used for YT yield accounting.

\paragraph{For ERC-4626 Vaults.}
\[
\texttt{exchange\_rate} = \texttt{vault.convertToAssets}(W)
\]

\paragraph{For Oracle-Based Tokens.}
\[
\texttt{exchange\_rate} = \texttt{oracle.get\_index}()
\]

\paragraph{Invariant.} The exchange rate is monotonically non-decreasing due to watermark protection. This is enforced by the global index update logic, not by the SY contract itself.

\section{PT/YT Minting Mechanism}

\subsection{Minting Authority Invariant}

Only the YT contract may mint or burn PT and YT tokens. This invariant is enforced by access control on the \texttt{mint} and \texttt{burn} functions of both token contracts.

\begin{center}
\begin{tabular}{lll}
\toprule
\textbf{Token} & \textbf{Authorized Caller} & \textbf{Enforcement} \\
\midrule
PT & YT contract only & $\texttt{require}(caller = \texttt{yt\_address})$ \\
YT & YT contract only & Internal function (self-mint) \\
\bottomrule
\end{tabular}
\end{center}

\subsection{Index Update Triggers}

The global index and user interest are updated at specific points to ensure correct yield accounting.

\begin{center}
\begin{tabular}{lcc}
\toprule
\textbf{Function} & \textbf{Updates Global Index} & \textbf{Updates User Interest} \\
\midrule
\texttt{mint\_py} & Yes & Yes (receiver) \\
\texttt{redeem\_py} & Yes & Yes (caller) \\
\texttt{redeem\_pt\_after\_expiry} & No & No \\
\texttt{claim\_interest} & Yes & Yes (caller) \\
\texttt{YT.transfer} & Yes & Yes (sender and recipient) \\
\texttt{YT.transferFrom} & Yes & Yes (from and to) \\
\bottomrule
\end{tabular}
\end{center}

\paragraph{YT Transfer Hook.} YT transfers must update both sender and recipient interest before modifying balances. Failure to do so would cause interest accounting errors. This is implemented via an internal \texttt{\_before\_token\_transfer} hook.

\subsection{Mint PY (Split SY to PT + YT)}

\begin{algorithm}[H]
\caption{Mint Principal and Yield Tokens}
\underline{function \textsc{MintPY}}$(receiver, amount_{sy}) \rightarrow (pt, yt)$\;
\Input{$receiver$: address, $amount_{sy}$: u256}
\Output{$(pt, yt)$: amounts of PT and YT minted}
\Preconditions{
  $\texttt{block.timestamp} < expiry$\;
  $receiver \neq 0$\;
  $amount_{sy} > 0$\;
  $\texttt{SY.allowance}(caller, YT) \geq amount_{sy}$\;
}
\BlankLine
\tcp{Update indices before balance changes}
$I_{global} \gets \max(I_{global}, \texttt{current\_exchange\_rate})$\;
$\textsc{UpdateUserInterest}(receiver)$\;
\BlankLine
\tcp{Transfer SY from caller to YT contract}
$\texttt{SY.balanceOf}(caller) \gets \texttt{SY.balanceOf}(caller) - amount_{sy}$\;
$\texttt{SY.balanceOf}(YT) \gets \texttt{SY.balanceOf}(YT) + amount_{sy}$\;
\BlankLine
\tcp{Mint equal PT and YT to receiver}
$\texttt{PT.balanceOf}(receiver) \gets \texttt{PT.balanceOf}(receiver) + amount_{sy}$\;
$\texttt{PT.totalSupply} \gets \texttt{PT.totalSupply} + amount_{sy}$\;
$\texttt{YT.balanceOf}(receiver) \gets \texttt{YT.balanceOf}(receiver) + amount_{sy}$\;
$\texttt{YT.totalSupply} \gets \texttt{YT.totalSupply} + amount_{sy}$\;
\BlankLine
\textbf{emit} $\textsc{MintPY}(caller, receiver, amount_{sy}, amount_{sy})$\;
\Return $(amount_{sy}, amount_{sy})$\;
\BlankLine
\textbf{Post-invariant:} $\texttt{PT.totalSupply} = \texttt{YT.totalSupply}$\;
\BlankLine
\Reverts{
  $\texttt{block.timestamp} \geq expiry$: ``expired''\;
  $amount_{sy} = 0$: ``zero amount''\;
  $receiver = 0$: ``zero address''\;
  allowance insufficient: ``ERC20: insufficient allowance''\;
}
\end{algorithm}

\subsection{Interest Accrual}

\begin{algorithm}[H]
\caption{Update User Interest}
\underline{function \textsc{UpdateUserInterest}}$(user)$\;
\Input{$user$: address}
\Preconditions{$I_{global}$ has been updated to current value}
\BlankLine
\eIf{$I_{user}[user] = 0$}{
  $I_{user}[user] \gets I_{global}$\tcp*{First interaction}
  \Return\;
}{
  $pending \gets \texttt{YT.balanceOf}(user) \times \dfrac{I_{global} - I_{user}[user]}{I_{user}[user]}$\;
  $\texttt{interest\_owed}[user] \gets \texttt{interest\_owed}[user] + pending$\;
  $I_{user}[user] \gets I_{global}$\;
}
\end{algorithm}

\section{Redemption Mechanisms}

\subsection{Pre-Expiry Redemption (PT + YT to SY)}

\begin{algorithm}[H]
\caption{Redeem PT and YT for SY (Pre-Expiry)}
\underline{function \textsc{RedeemPY}}$(receiver, amount) \rightarrow sy$\;
\Input{$receiver$: address, $amount$: u256}
\Output{$sy$: amount of SY returned}
\Preconditions{
  $\texttt{block.timestamp} < expiry$\;
  $receiver \neq 0$\;
  $amount > 0$\;
  $\texttt{PT.balanceOf}(caller) \geq amount$\;
  $\texttt{YT.balanceOf}(caller) \geq amount$\;
}
\BlankLine
\tcp{Update indices before balance changes}
$I_{global} \gets \max(I_{global}, \texttt{current\_exchange\_rate})$\;
$\textsc{UpdateUserInterest}(caller)$\;
\BlankLine
\tcp{Burn equal PT and YT from caller}
$\texttt{PT.balanceOf}(caller) \gets \texttt{PT.balanceOf}(caller) - amount$\;
$\texttt{PT.totalSupply} \gets \texttt{PT.totalSupply} - amount$\;
$\texttt{YT.balanceOf}(caller) \gets \texttt{YT.balanceOf}(caller) - amount$\;
$\texttt{YT.totalSupply} \gets \texttt{YT.totalSupply} - amount$\;
\BlankLine
\tcp{Transfer SY from YT contract to receiver}
$\texttt{SY.balanceOf}(YT) \gets \texttt{SY.balanceOf}(YT) - amount$\;
$\texttt{SY.balanceOf}(receiver) \gets \texttt{SY.balanceOf}(receiver) + amount$\;
\BlankLine
\textbf{emit} $\textsc{RedeemPY}(caller, receiver, amount)$\;
\Return $amount$\;
\BlankLine
\textbf{Post-invariant:} $\texttt{PT.totalSupply} = \texttt{YT.totalSupply}$\;
\BlankLine
\Reverts{
  $\texttt{block.timestamp} \geq expiry$: ``use redeem\_pt\_after\_expiry''\;
  $amount = 0$: ``zero amount''\;
  $receiver = 0$: ``zero address''\;
  PT balance insufficient: ``insufficient PT''\;
  YT balance insufficient: ``insufficient YT''\;
}
\end{algorithm}

\subsection{Post-Expiry Redemption (PT to SY)}

\begin{algorithm}[H]
\caption{Redeem PT for SY (Post-Expiry)}
\underline{function \textsc{RedeemPTAfterExpiry}}$(receiver, amount) \rightarrow sy$\;
\Input{$receiver$: address, $amount$: u256}
\Output{$sy$: amount of SY returned (guaranteed 1:1)}
\Preconditions{
  $\texttt{block.timestamp} \geq expiry$\;
  $receiver \neq 0$\;
  $amount > 0$\;
  $\texttt{PT.balanceOf}(caller) \geq amount$\;
}
\BlankLine
\tcp{No index update needed post-expiry}
$\texttt{PT.balanceOf}(caller) \gets \texttt{PT.balanceOf}(caller) - amount$\;
$\texttt{PT.totalSupply} \gets \texttt{PT.totalSupply} - amount$\;
$\texttt{SY.balanceOf}(YT) \gets \texttt{SY.balanceOf}(YT) - amount$\;
$\texttt{SY.balanceOf}(receiver) \gets \texttt{SY.balanceOf}(receiver) + amount$\;
\BlankLine
\textbf{emit} $\textsc{RedeemPTAfterExpiry}(caller, receiver, amount)$\;
\Return $amount$\;
\BlankLine
\Reverts{
  $\texttt{block.timestamp} < expiry$: ``not expired''\;
  $amount = 0$: ``zero amount''\;
  $receiver = 0$: ``zero address''\;
  PT balance insufficient: ``insufficient PT''\;
}
\end{algorithm}

\subsection{YT Interest Claim}

\begin{algorithm}[H]
\caption{Claim YT Interest}
\underline{function \textsc{ClaimInterest}}$(receiver) \rightarrow interest$\;
\Input{$receiver$: address}
\Output{$interest$: amount of SY claimed}
\Preconditions{$receiver \neq 0$}
\BlankLine
$I_{global} \gets \max(I_{global}, \texttt{current\_exchange\_rate})$\;
$\textsc{UpdateUserInterest}(caller)$\;
\BlankLine
$interest \gets \texttt{interest\_owed}[caller]$\;
$\texttt{interest\_owed}[caller] \gets 0$\;
\BlankLine
$\texttt{SY.balanceOf}(YT) \gets \texttt{SY.balanceOf}(YT) - interest$\;
$\texttt{SY.balanceOf}(receiver) \gets \texttt{SY.balanceOf}(receiver) + interest$\;
\BlankLine
\textbf{emit} $\textsc{ClaimInterest}(caller, receiver, interest)$\;
\Return $interest$\;
\BlankLine
\Reverts{
  $receiver = 0$: ``zero address''\;
  $\texttt{SY.balanceOf}(YT) < interest$: ``insufficient SY reserves''\;
}
\end{algorithm}

\section{AMM Mechanism}

The Market contract implements a logit-based AMM for PT/SY trading with time-decay characteristics that ensure PT price converges to 1 at maturity.

\subsection{Implied Rate Storage}

The market stores the last implied rate to maintain price continuity across trades.

\begin{center}
\begin{tabular}{llll}
\toprule
\textbf{Variable} & \textbf{Type} & \textbf{Units} & \textbf{Description} \\
\midrule
$r_{last}$ & u256 & WAD-scaled & Natural log of implied annualized rate \\
\bottomrule
\end{tabular}
\end{center}

\paragraph{Update formula.} After each trade, the implied rate is recalculated from the new pool state:
\[
r_{new} = \ln(\text{exchangeRate}) \times \frac{Y}{\tau}
\]
where $Y = 31{,}536{,}000$ (seconds per year) and $\tau$ is time to expiry in seconds. The value is clamped to $[0, r_{max}]$ where $r_{max} = 4.6 \times 10^{18}$.

\subsection{Precompute Values}

\begin{algorithm}[H]
\caption{Precompute Market Values}
\underline{function \textsc{Precompute}}$(state, t_{now}) \rightarrow pre$\;
\Input{$state$: MarketState, $t_{now}$: u64 (current timestamp)}
\Output{$pre$: MarketPrecompute struct}
\Preconditions{$t_{now} < state.expiry$}
\BlankLine
$\tau \gets state.expiry - t_{now}$\;
$r_{scalar} \gets \dfrac{state.scalar\_root \times Y}{\tau}$\;
$p \gets \dfrac{state.S_{PT}}{state.S_{PT} + state.S_{SY}}$\;
\BlankLine
$\text{exchangeRate}_{target} \gets \exp\left(\dfrac{r_{last} \times \tau}{Y}\right)$\;
$r_{anchor} \gets \text{exchangeRate}_{target} - \dfrac{\operatorname{logit}(p)}{r_{scalar}}$\;
\BlankLine
$\phi_{adjusted} \gets \exp\left(\dfrac{\ln(\phi_{root}) \times \tau}{Y}\right)$\;
\BlankLine
\Return $\{r_{scalar}, r_{anchor}, \phi_{adjusted}, \tau\}$\;
\BlankLine
\Reverts{
  $\tau = 0$: ``expired''\;
  $p \notin [p_{min}, p_{max}]$: ``invalid state''\;
}
\end{algorithm}

\subsection{Price Calculation}

The exchange rate depends on the current pool proportion:
\[
\text{exchangeRate} = \frac{\operatorname{logit}(p)}{r_{scalar}} + r_{anchor}
\]
where $p = \frac{S_{PT}}{S_{PT} + S_{SY}}$ is the proportion of PT in the pool.

As expiry approaches, $r_{scalar} \to \infty$, which flattens the curve and forces the exchange rate toward 1.

\section{Swap Mechanics}

All swap functions reject trades that would violate proportion bounds or occur after expiry.

\subsection{Swap Exact PT for SY}

\begin{algorithm}[H]
\caption{Swap Exact PT for SY}
\underline{function \textsc{SwapExactPTForSY}}$(receiver, pt_{in}, sy_{min}) \rightarrow sy_{out}$\;
\Input{$receiver$: address, $pt_{in}$: u256, $sy_{min}$: u256 (slippage bound)}
\Output{$sy_{out}$: amount of SY received}
\Preconditions{
  $\texttt{block.timestamp} < expiry$\;
  $receiver \neq 0$\;
  $pt_{in} > 0$\;
  $\texttt{PT.allowance}(caller, Market) \geq pt_{in}$\;
}
\BlankLine
$pre \gets \textsc{Precompute}(state, \texttt{block.timestamp})$\;
\BlankLine
\tcp{Calculate new proportion after adding PT}
$S'_{PT} \gets state.S_{PT} + pt_{in}$\;
$p' \gets \dfrac{S'_{PT}}{S'_{PT} + state.S_{SY}}$\;
\lIf{$p' > p_{max}$}{\Revert ``proportion bounds exceeded''}
\BlankLine
\tcp{Calculate exchange rate and output}
$\text{exchangeRate} \gets \dfrac{\operatorname{logit}(p')}{pre.r_{scalar}} + pre.r_{anchor}$\;
\lIf{$\text{exchangeRate} \leq 0$}{\Revert ``invalid exchange rate''}
\BlankLine
$sy_{gross} \gets \dfrac{pt_{in} \times \text{exchangeRate}}{W}$\;
$fee \gets \dfrac{sy_{gross} \times (W - pre.\phi_{adjusted})}{W}$\;
$sy_{out} \gets sy_{gross} - fee$\;
\BlankLine
\lIf{$sy_{out} < sy_{min}$}{\Revert ``slippage exceeded''}
\lIf{$state.S_{SY} < sy_{out}$}{\Revert ``insufficient reserves''}
\BlankLine
\tcp{State changes}
$state.S_{PT} \gets state.S_{PT} + pt_{in}$\;
$state.S_{SY} \gets state.S_{SY} - sy_{out}$\;
$r_{last} \gets \text{clamp}\left(\ln(\text{exchangeRate}) \times \dfrac{Y}{\tau}, 0, r_{max}\right)$\;
\BlankLine
\tcp{Token transfers}
$\texttt{PT.transferFrom}(caller, Market, pt_{in})$\;
$\texttt{SY.transfer}(receiver, sy_{out})$\;
\BlankLine
\textbf{emit} $\textsc{Swap}(caller, receiver, pt_{in}, sy_{out})$\;
\Return $sy_{out}$\;
\end{algorithm}

\subsection{Swap Exact SY for PT}

\begin{algorithm}[H]
\caption{Swap Exact SY for PT}
\underline{function \textsc{SwapExactSYForPT}}$(receiver, sy_{in}, pt_{min}) \rightarrow pt_{out}$\;
\Input{$receiver$: address, $sy_{in}$: u256, $pt_{min}$: u256 (slippage bound)}
\Output{$pt_{out}$: amount of PT received}
\Preconditions{
  $\texttt{block.timestamp} < expiry$\;
  $receiver \neq 0$\;
  $sy_{in} > 0$\;
  $\texttt{SY.allowance}(caller, Market) \geq sy_{in}$\;
}
\BlankLine
$pre \gets \textsc{Precompute}(state, \texttt{block.timestamp})$\;
\BlankLine
\tcp{Apply fee to input}
$fee \gets \dfrac{sy_{in} \times (W - pre.\phi_{adjusted})}{W}$\;
$sy_{net} \gets sy_{in} - fee$\;
$S'_{SY} \gets state.S_{SY} + sy_{net}$\;
\BlankLine
\tcp{Solve for $pt_{out}$ via bisection}
$lo \gets 0$\;
$hi \gets state.S_{PT} - L_{min}$\;
\For{$i \gets 1$ \KwTo $64$}{
  $mid \gets (lo + hi) / 2$\;
  $p' \gets \dfrac{state.S_{PT} - mid}{state.S_{PT} - mid + S'_{SY}}$\;
  $\text{rate} \gets \dfrac{\operatorname{logit}(p')}{pre.r_{scalar}} + pre.r_{anchor}$\;
  $\text{implied} \gets \dfrac{sy_{net} \times W}{\text{rate}}$\;
  \eIf{$\text{implied} > mid$}{$lo \gets mid$}{$hi \gets mid$}
  \lIf{$|hi - lo| \leq 1$}{\textbf{break}}
}
$pt_{out} \gets lo$\;
\BlankLine
$p' \gets \dfrac{state.S_{PT} - pt_{out}}{state.S_{PT} - pt_{out} + S'_{SY}}$\;
\lIf{$p' < p_{min}$}{\Revert ``proportion bounds exceeded''}
\lIf{$pt_{out} < pt_{min}$}{\Revert ``slippage exceeded''}
\BlankLine
\tcp{State changes and transfers}
$state.S_{SY} \gets state.S_{SY} + sy_{net}$\;
$state.S_{PT} \gets state.S_{PT} - pt_{out}$\;
Update $r_{last}$ as in \textsc{SwapExactPTForSY}\;
\BlankLine
$\texttt{SY.transferFrom}(caller, Market, sy_{in})$\;
$\texttt{PT.transfer}(receiver, pt_{out})$\;
\BlankLine
\textbf{emit} $\textsc{Swap}(caller, receiver, sy_{in}, pt_{out})$\;
\Return $pt_{out}$\;
\end{algorithm}

\paragraph{Bisection rationale.} The SY-to-PT direction requires solving for $pt_{out}$ implicitly because the exchange rate depends on the post-trade proportion. Bisection is monotonic, bounded, and guaranteed to converge within 64 iterations for WAD-scaled values.

\subsection{Inverse Swaps}

\begin{algorithm}[H]
\caption{Swap SY for Exact PT}
\underline{function \textsc{SwapSYForExactPT}}$(receiver, pt_{out}, sy_{max}) \rightarrow sy_{spent}$\;
Inverse of \textsc{SwapExactSYForPT}: given $pt_{out}$, solve for $sy_{in}$.\;
\Reverts{$sy_{spent} > sy_{max}$: ``slippage exceeded''}
\end{algorithm}

\begin{algorithm}[H]
\caption{Swap PT for Exact SY}
\underline{function \textsc{SwapPTForExactSY}}$(receiver, sy_{out}, pt_{max}) \rightarrow pt_{spent}$\;
Inverse of \textsc{SwapExactPTForSY}: given $sy_{out}$, solve for $pt_{in}$.\;
\Reverts{$pt_{spent} > pt_{max}$: ``slippage exceeded''}
\end{algorithm}

\section{Liquidity Provision}

\subsection{Add Liquidity (Mint LP)}

\begin{algorithm}[H]
\caption{Add Liquidity}
\underline{function \textsc{Mint}}$(receiver, sy_{desired}, pt_{desired}) \rightarrow (sy, pt, lp)$\;
\Input{$receiver$: address, $sy_{desired}$: u256, $pt_{desired}$: u256}
\Output{$(sy, pt, lp)$: amounts used and LP tokens minted}
\Preconditions{
  $\texttt{block.timestamp} < expiry$\;
  $receiver \neq 0$\;
  $sy_{desired} > 0 \land pt_{desired} > 0$\;
}
\BlankLine
\eIf{$L_{total} = 0$}{
  \tcp{Initial liquidity}
  $lp \gets \sqrt{sy_{desired} \times pt_{desired}} - L_{min}$\;
  $sy \gets sy_{desired}$\;
  $pt \gets pt_{desired}$\;
  \tcp{Lock minimum to dead address}
  $\textsc{MintLP}(\texttt{DEAD\_ADDRESS}, L_{min})$\;
}{
  \tcp{Proportional liquidity}
  $r_{sy} \gets \dfrac{sy_{desired} \times W}{S_{SY}}$\;
  $r_{pt} \gets \dfrac{pt_{desired} \times W}{S_{PT}}$\;
  $r \gets \min(r_{sy}, r_{pt})$\;
  \BlankLine
  $sy \gets \dfrac{r \times S_{SY}}{W}$\;
  $pt \gets \dfrac{r \times S_{PT}}{W}$\;
  $lp \gets \dfrac{r \times L_{total}}{W}$\;
}
\BlankLine
\tcp{Verify proportion bounds after adding}
$p' \gets \dfrac{S_{PT} + pt}{S_{PT} + pt + S_{SY} + sy}$\;
\lIf{$p' \notin [p_{min}, p_{max}]$}{\Revert ``proportion bounds''}
\BlankLine
$S_{SY} \gets S_{SY} + sy$\;
$S_{PT} \gets S_{PT} + pt$\;
$L_{total} \gets L_{total} + lp$\;
\BlankLine
\textbf{emit} $\textsc{Mint}(caller, receiver, sy, pt, lp)$\;
\Return $(sy, pt, lp)$\;
\BlankLine
\Reverts{
  expired, zero address, zero amounts, $lp = 0$, proportion bounds\;
}
\end{algorithm}

\paragraph{Constants.}
$L_{min} = 1000$ (locked to dead address on first mint). $\texttt{DEAD\_ADDRESS} = \texttt{0x...001}$.

\subsection{Remove Liquidity (Burn LP)}

\begin{algorithm}[H]
\caption{Remove Liquidity}
\underline{function \textsc{Burn}}$(receiver, lp) \rightarrow (sy_{out}, pt_{out})$\;
\Input{$receiver$: address, $lp$: u256}
\Output{$(sy_{out}, pt_{out})$: amounts received}
\Preconditions{
  $receiver \neq 0$\;
  $lp > 0$\;
  $\texttt{LP.balanceOf}(caller) \geq lp$\;
  $L_{total} - lp \geq L_{min}$\;
}
\BlankLine
$r \gets \dfrac{lp \times W}{L_{total}}$\;
$sy_{out} \gets \dfrac{r \times S_{SY}}{W}$\;
$pt_{out} \gets \dfrac{r \times S_{PT}}{W}$\;
\BlankLine
$L_{total} \gets L_{total} - lp$\;
$S_{SY} \gets S_{SY} - sy_{out}$\;
$S_{PT} \gets S_{PT} - pt_{out}$\;
\BlankLine
$\texttt{LP.burn}(caller, lp)$\;
$\texttt{SY.transfer}(receiver, sy_{out})$\;
$\texttt{PT.transfer}(receiver, pt_{out})$\;
\BlankLine
\textbf{emit} $\textsc{Burn}(caller, receiver, sy_{out}, pt_{out}, lp)$\;
\Return $(sy_{out}, pt_{out})$\;
\BlankLine
\textbf{Note:} Burning is allowed after expiry. LPs can always exit.\;
\end{algorithm}

\section{Key Invariants}

The protocol maintains two categories of invariants: safety invariants enforced by code, and economic invariants maintained by market incentives.

\subsection{Safety Invariants (Enforced by Code)}

These invariants are checked programmatically and violations cause transaction reversion.

\begin{center}
\begin{tabularx}{\textwidth}{l >{\raggedright\arraybackslash}X l}
\toprule
\textbf{Invariant} & \textbf{Expression} & \textbf{Enforcement} \\
\midrule
Supply equality & $\text{PT.totalSupply} = \text{YT.totalSupply}$ & Atomic mint/burn in YT contract \\
Index watermark & $I_{global}^{t+1} \geq I_{global}^{t}$ & $\max()$ in update function \\
Reserve positivity & $S_{SY} > 0 \land S_{PT} > 0$ & $L_{min}$ lock \\
Proportion bounds & $p_{min} \leq p \leq p_{max}$ & Trade rejection \\
Rate bounds & $0 \leq r \leq r_{max}$ & Clamping before storage \\
PT/YT mint authority & Only YT contract mints/burns & Access control check \\
\bottomrule
\end{tabularx}
\end{center}

\subsection{Economic Invariants (Maintained by Arbitrage)}

These invariants are expected to hold due to economic incentives but are not enforced on-chain.

\begin{center}
\begin{tabularx}{\textwidth}{l >{\raggedright\arraybackslash}X l}
\toprule
\textbf{Invariant} & \textbf{Expression} & \textbf{Mechanism} \\
\midrule
Conservation & $P_{PT} + P_{YT} \approx 1$ (in SY terms) & Arbitrage between mint/redeem and market \\
PT price convergence & $P_{PT} \to 1$ as $\tau \to 0$ & AMM curve design + arbitrage \\
Implied rate consistency & Market rate $\approx$ fair rate & Arbitrage against external markets \\
\bottomrule
\end{tabularx}
\end{center}

\paragraph{Violation consequences.}
\begin{itemize}
\item Supply equality violation: Accounting mismatch; claims exceed or fall short of reserves.
\item Conservation deviation: Temporary mispricing; corrected by arbitrage.
\item Index watermark violation: Loss of accrued yield for YT holders.
\end{itemize}

\section{Failure Modes and Recovery}

Each failure mode is documented with detection mechanism, on-chain mitigation, and residual risk.

\subsection{Oracle Failure}

\begin{center}
\begin{tabularx}{\textwidth}{l >{\raggedright\arraybackslash}X}
\toprule
\textbf{Aspect} & \textbf{Description} \\
\midrule
Detection & Oracle returns stale data (age $>$ MAX\_ORACLE\_AGE, default 24 hours) \\
On-chain response & Watermark index used; new yield not distributed until recovery \\
Trigger & Permissionless (watermark logic is automatic) \\
Residual risk & Yield generated during outage may not accrue to YT holders \\
\bottomrule
\end{tabularx}
\end{center}

\subsection{Underlying Depeg}

\begin{center}
\begin{tabularx}{\textwidth}{l >{\raggedright\arraybackslash}X}
\toprule
\textbf{Scenario} & \textbf{Behavior} \\
\midrule
Temporary depeg & Index stays at watermark; no new yield distributed. Auto-recovery when rate exceeds watermark. \\
Permanent depeg & YT stops accruing yield indefinitely. PT redemption unaffected if SY remains solvent. \\
Complete failure & SY loses value. Users should exit via redemption. No protocol-level recovery. \\
\bottomrule
\end{tabularx}
\end{center}

\subsection{Pool Imbalance}

\begin{center}
\begin{tabularx}{\textwidth}{l >{\raggedright\arraybackslash}X}
\toprule
\textbf{Aspect} & \textbf{Description} \\
\midrule
Detection & Proportion approaches $p_{min}$ or $p_{max}$ \\
On-chain response & Trades that exceed bounds are rejected \\
Self-correction & Extreme proportions create large arbitrage opportunities \\
Residual risk & Temporarily reduced liquidity in one direction \\
\bottomrule
\end{tabularx}
\end{center}

\subsection{Flash Loan Attacks}

\begin{center}
\begin{tabularx}{\textwidth}{l >{\raggedright\arraybackslash}X}
\toprule
\textbf{Attack Vector} & \textbf{Mitigation} \\
\midrule
Price manipulation & Logit curve provides increasing resistance at proportion extremes \\
Sandwich attacks & User-specified slippage protection ($sy_{min}$ / $pt_{max}$) \\
Oracle manipulation & Pragma TWAP with configurable window (default: 1 hour) \\
\bottomrule
\end{tabularx}
\end{center}

\subsection{Expired Market Operations}

\begin{center}
\begin{tabular}{lll}
\toprule
\textbf{Operation} & \textbf{Before Expiry} & \textbf{After Expiry} \\
\midrule
Swap & Allowed & Rejected \\
Mint LP & Allowed & Rejected \\
Burn LP & Allowed & Allowed \\
Redeem PT + YT & Allowed & Use \textsc{RedeemPTAfterExpiry} \\
Redeem PT only & Rejected & Allowed (1:1) \\
Claim Interest & Allowed & Allowed \\
\bottomrule
\end{tabular}
\end{center}


\chapter{Economic Analysis \& Incentives}

This section analyzes the economic properties of the protocol, including actor incentives, equilibrium conditions, and behavior under stress.

\paragraph{Scope and assumptions.} All results in this section are conditional on the external assumptions stated in Chapter 3: oracle liveness, underlying asset solvency, and bounded transaction costs. Claims about equilibrium behavior assume competitive arbitrage and sufficient market liquidity.

\section{Actors and Incentives}

\subsection{PT Holders (Fixed-Rate Exposure)}

\paragraph{Decision rule.}
A rational actor buys PT when the market-implied fixed rate exceeds their required return:
\[
r_{\mathrm{implied}} \geq r_{\mathrm{alt}} + \pi
\]
where $r_{\mathrm{alt}}$ is the opportunity cost of capital and $\pi$ is a risk premium covering contract risk, oracle risk, and underlying impairment risk.

\paragraph{Implied rate.}
Let $\Delta$ be time-to-expiry in years and $P_{PT}$ the PT price in SY units. The implied annualized return is:
\[
r_{\mathrm{implied}} = P_{PT}^{-1/\Delta} - 1
\]
For example, if $P_{PT} = 0.95$ and $\Delta = 0.5$ years, the implied APY is $(0.95)^{-2} - 1 \approx 10.8\%$.

\paragraph{Risks.}
\begin{itemize}
\item \textbf{Smart contract risk}: Low probability, high severity. Mitigated by audits and formal verification.
\item \textbf{Underlying impairment}: PT redemption is deterministic in SY units, but SY value depends on underlying solvency.
\item \textbf{Opportunity cost}: Capital is locked; holder cannot chase higher yields elsewhere.
\item \textbf{Early exit}: Must sell on AMM at prevailing price, subject to slippage and mark-to-market loss.
\end{itemize}

\paragraph{Implications.}
PT provides deterministic redemption in SY units at expiry, conditional on SY solvency. Prior to expiry, PT is mark-to-market and may trade at a discount or premium depending on rate expectations and liquidity.

\subsection{YT Holders (Yield Speculation)}

\paragraph{Decision rule.}
A rational actor buys YT when expected yield exceeds the YT price:
\[
\mathbb{E}[g(\Delta)] \geq P_{YT}
\]
where $g(\Delta)$ is the expected fractional yield growth of the underlying over time-to-expiry $\Delta$. Equivalently, if $y$ is the expected APY:
\[
1 - e^{-y\Delta} \geq P_{YT}
\]
This directly ties to the conservation identity $P_{PT} + P_{YT} = 1$.

\paragraph{Leverage effect.}
YT provides leveraged exposure to yield. If $P_{YT} = 0.05$ (5\% of SY value) and the underlying yields 10\% APY over a 6-month period:
\begin{itemize}
\item SY holder earns: $\approx 5\%$ on 100\% capital
\item YT holder earns: $\approx 5\%$ yield on 5\% capital = $\approx 100\%$ effective return on YT investment
\end{itemize}

\paragraph{Risks.}
\begin{itemize}
\item \textbf{Time decay}: YT value converges to zero as expiry approaches, regardless of yield.
\item \textbf{Yield drops}: Reduced interest collection; YT loses value rapidly.
\item \textbf{Complete depeg}: If underlying stops generating yield, YT becomes worthless.
\end{itemize}

\paragraph{Implications.}
YT is a high-risk, high-reward instrument suitable for actors with strong convictions about future yield. Longer maturities provide more leverage but greater time decay exposure.

\subsection{Liquidity Providers}

\paragraph{Decision rule.}
LP provision is attractive when expected return exceeds opportunity cost:
\[
\mathbb{E}[\text{Fee APY}] + r_{SY} \cdot w_{SY} - \mathbb{E}[\text{IL}] \geq r_{\mathrm{alt}}
\]
where $r_{SY}$ is the SY yield rate, $w_{SY}$ is the SY weight in the pool, and IL is impermanent loss.

\paragraph{Impermanent loss dynamics.}
Unlike traditional AMMs where IL is symmetric, PT/SY pools exhibit asymmetric and time-dependent IL:

\begin{center}
\begin{tabular}{lll}
\toprule
\textbf{Rate Movement} & \textbf{PT Price} & \textbf{LP Impact} \\
\midrule
Rates increase & PT decreases & IL: LP holds excess PT \\
Rates decrease & PT increases & IL: LP holds excess SY \\
Near expiry & PT $\to$ 1 & IL minimizes as prices converge \\
\bottomrule
\end{tabular}
\end{center}

\paragraph{Time-to-expiry effect.}
As expiry approaches, the PT price becomes less sensitive to rate changes (via the rate scalar $r_{scalar} \to \infty$), reducing volatility-driven IL. However, this also typically reduces trading volume and fee opportunities. LPs face a tradeoff between IL protection and fee income as maturity nears.

\paragraph{Risks.}
\begin{itemize}
\item \textbf{Directional IL}: Unlike xy=k pools, rate movements cause directional exposure shifts.
\item \textbf{Rate volatility}: High volatility increases IL without proportional fee compensation.
\item \textbf{Low volume periods}: Fee income may not offset capital opportunity cost.
\end{itemize}

\paragraph{Implications.}
LP returns are path-dependent and sensitive to rate volatility. Optimal entry is during stable rate environments; exit timing around volatility events is critical.

\subsection{Arbitrageurs}

\paragraph{Conservation bounds.}
The conservation identity creates soft bounds on prices:
\[
P_{PT} + P_{YT} = 1 \pm \epsilon
\]
where $\epsilon$ is bounded by transaction costs, slippage, and execution latency. In liquid markets with low gas costs, $\epsilon < 0.5\%$ is typical.

\paragraph{Arbitrage mechanisms.}

\begin{center}
\begin{tabular}{lll}
\toprule
\textbf{Condition} & \textbf{Action} & \textbf{Effect} \\
\midrule
$P_{PT} + P_{YT} > 1 + \epsilon$ & Mint PT+YT, sell both & Prices decrease \\
$P_{PT} + P_{YT} < 1 - \epsilon$ & Buy PT+YT, redeem for SY & Prices increase \\
Horizon yield $>$ external & Sell PT (or buy YT) & Implied yield decreases \\
Horizon yield $<$ external & Buy PT (or sell YT) & Implied yield increases \\
\bottomrule
\end{tabular}
\end{center}

\paragraph{Implications.}
Arbitrage activity maintains price consistency with external rate expectations and enforces the conservation identity within transaction cost bounds. Arbitrage profitability decreases as markets become more efficient.

\section{Equilibrium Conditions}

\subsection{Price Equilibrium}

At equilibrium, the implied yield reflects market consensus on future rates:
\[
r_{\mathrm{implied}} = \mathbb{E}[r_{\mathrm{actual}}] + \pi
\]
where $\pi$ is a risk premium for contract and oracle risk. Deviations from this condition create arbitrage opportunities that are corrected by rational traders.

\subsection{Liquidity Equilibrium}

LP provision reaches equilibrium when marginal return equals opportunity cost:
\[
\text{Fee APY} + r_{SY} \cdot w_{SY} - \mathbb{E}[\text{IL}] = r_{\mathrm{alt}}
\]

Factors affecting LP equilibrium depth:
\begin{itemize}
\item Higher trading volume $\Rightarrow$ more fees $\Rightarrow$ increased LP supply
\item Higher rate volatility $\Rightarrow$ more IL $\Rightarrow$ decreased LP supply
\item Higher underlying yield $\Rightarrow$ more attractive SY exposure $\Rightarrow$ increased LP supply
\end{itemize}

\section{Stress Scenarios}

\subsection{Yield Collapse}

\paragraph{Scenario.} Underlying yield drops from 10\% to 2\% APY.

\begin{center}
\begin{tabular}{lll}
\toprule
\textbf{Actor} & \textbf{Immediate Impact} & \textbf{Final Outcome} \\
\midrule
PT Holder & Mark-to-market decline & Converges to 1 SY at maturity \\
YT Holder & YT price drops 80\%+ & Collects only 2\% yield \\
LP & IL + reduced fee volume & Depends on exit timing \\
\bottomrule
\end{tabular}
\end{center}

\paragraph{Key observation.} PT holders who hold to maturity are unaffected in SY terms; the ``loss'' is opportunity cost relative to new higher-yielding alternatives.

\subsection{Yield Spike}

\paragraph{Scenario.} Underlying yield spikes from 10\% to 25\% APY.

\begin{center}
\begin{tabular}{lll}
\toprule
\textbf{Actor} & \textbf{Immediate Impact} & \textbf{Final Outcome} \\
\midrule
PT Holder & Opportunity cost (missed yield) & Still receives 1 SY at maturity \\
YT Holder & Large windfall & Collects 25\% yield; YT appreciates \\
LP & Directional IL & Higher volume partially offsets \\
\bottomrule
\end{tabular}
\end{center}

\subsection{Mass Redemption}

\paragraph{Scenario.} 80\% of PT+YT holders redeem simultaneously.

\paragraph{Key property.} Unlike lending protocols, the protocol has no maturity mismatch or run-induced insolvency risk, assuming SY is fully collateralized by the underlying. Each PT+YT pair is backed by exactly 1 SY held in the YT contract; redemption is always possible.

\paragraph{Constraints.}
\begin{itemize}
\item Redemption throughput is limited by on-chain block capacity.
\item AMM exits are subject to price impact and slippage.
\item Solvency does not depend on borrower repayment (unlike lending protocols).
\end{itemize}

\subsection{Underlying Depeg}

\paragraph{Scenario.} Underlying token loses 30\% of external value.

\begin{center}
\begin{tabular}{ll}
\toprule
\textbf{Asset} & \textbf{Impact} \\
\midrule
SY & Worth 0.7$\times$ previous external value \\
PT & Redeems for 1 SY (now worth 0.7$\times$) \\
YT & Yield continues at depressed rate \\
LP & Proportional loss on both reserve assets \\
\bottomrule
\end{tabular}
\end{center}

\paragraph{Key distinction.} Depeg affects the value of SY in external terms (e.g., USD), not the internal accounting identity PT $\to$ 1 SY. The protocol's internal guarantees remain intact; external value loss is passed through to all token holders.

\paragraph{Protocol response.}
\begin{itemize}
\item Watermark index prevents distributing negative yield.
\item All operations continue functioning normally.
\item Users can exit at prevailing market prices.
\end{itemize}

\section{Mechanism Properties}

\subsection{Equilibrium Characterization}

The system admits a candidate equilibrium where no actor can improve their expected payoff by unilaterally deviating:

\begin{center}
\begin{tabular}{ll}
\toprule
\textbf{Actor} & \textbf{Equilibrium Condition} \\
\midrule
PT Buyer & Implied yield = required return for risk \\
YT Buyer & Expected yield $\geq$ YT price + speculation premium \\
LP & Fee income = IL + opportunity cost \\
Arbitrageur & Spread $\leq$ execution costs \\
\bottomrule
\end{tabular}
\end{center}

Under competitive arbitrage and sufficient liquidity, deviations from equilibrium prices are mean-reverting.

\subsection{Incentive Properties}

\paragraph{Incentive compatibility.}
Honest behavior (accurate price discovery through trading) is locally incentive-compatible under competitive arbitrage and bounded manipulation costs. Price manipulation is costly (requires capital) and temporary (arbitrage corrects).

\paragraph{Individual rationality.}
Participation is voluntary. Expected utility is non-negative for informed participants who correctly assess risks and opportunities.

\paragraph{Budget balance.}
The protocol is self-sustaining: fees cover operational costs (gas for oracle updates, etc.). No external subsidies are required for equilibrium operation.

\section{Fee Economics}

\subsection{Fee Retention Model}

Fees are retained in pool reserves rather than extracted:
\begin{itemize}
\item Swap fees increase reserve balances (SY or PT depending on trade direction).
\item LP token value appreciates proportionally to accumulated fees.
\item Fee value is realized when LPs withdraw liquidity.
\end{itemize}

\paragraph{Terminology note.} The swap mechanism uses a \texttt{fee\_multiplier} $\in [0,1]$ where:
\[
\text{fee} = \text{gross\_amount} \times (1 - \text{fee\_multiplier})
\]
A fee multiplier of 0.999 corresponds to a 0.1\% fee rate.

\subsection{Optimal Fee Rate}

The optimal fee rate balances trader demand and LP incentives:
\begin{itemize}
\item \textbf{Too low}: Insufficient LP incentive $\Rightarrow$ thin liquidity $\Rightarrow$ high slippage.
\item \textbf{Too high}: Traders use alternatives $\Rightarrow$ low volume $\Rightarrow$ reduced LP returns.
\end{itemize}

Market forces determine the effective optimum:
\begin{itemize}
\item Competition between markets for the same underlying sets fee ceiling.
\item Arbitrage activity requires fees below arbitrage profit margins.
\item LP entry/exit adjusts liquidity depth to equilibrium levels.
\end{itemize}

\section{Value Accrual}

\subsection{Value Flow}

\begin{center}
\begin{tabular}{ll}
\toprule
\textbf{Value Source} & \textbf{Recipient} \\
\midrule
Underlying yield & YT holders (until expiry) \\
PT discount capture & PT holders (at maturity) \\
Swap fees & LP holders \\
Arbitrage spreads & Arbitrageurs \\
\bottomrule
\end{tabular}
\end{center}

\paragraph{Zero-sum framing.}
Discount capture (PT) and yield capture (YT) are two sides of the same decomposition: PT holders forgo yield in exchange for certainty; YT holders accept uncertainty for leveraged yield exposure. Fees and arbitrage spreads are paid by traders through execution price.

\subsection{Protocol Value Capture}

\paragraph{Current design.} Zero protocol fee. All value flows to users (PT/YT/LP holders and arbitrageurs).

\paragraph{Future consideration.} A governance-controlled protocol fee on swaps could direct value to a treasury or token holders. This would be implemented as a fraction of swap fees retained by the protocol rather than accruing to LPs.



\chapter{Risk Analysis \& Threat Model}

This section enumerates the risks associated with the protocol, the controls in place, and residual risks that users must accept. Each risk category follows a consistent schema: threat, impact, preconditions, controls, and residual risk.

\paragraph{Threat model assumptions.}
\begin{itemize}
\item The Starknet network operates correctly (liveness, validity proofs, L1 finality).
\item The Pragma oracle infrastructure is not fundamentally compromised.
\item Cryptographic primitives (hashing, signatures) remain secure.
\item Users can monitor on-chain events and act within timelock windows.
\end{itemize}

\section{Smart Contract Risk}

\subsection{Code Vulnerabilities}

\begin{center}
\begin{tabularx}{\textwidth}{l >{\raggedright\arraybackslash}X l l}
\toprule
\textbf{Threat} & \textbf{Controls} & \textbf{Status} & \textbf{Residual} \\
\midrule
Reentrancy & Checks-effects-interactions; explicit reentrancy guards; minimize external calls in state-mutating paths & Implemented & Low \\
Arithmetic overflow & Widened-precision math; explicit overflow checks; fuzz testing over boundary values & Implemented & Low \\
Access control bypass & Role-based access control (RBAC); explicit permission checks; separation of duties & Implemented & Low \\
Logic errors & Unit tests; integration tests; invariant fuzzing; external audit & Partial & Medium \\
Front-running / MEV & User-defined slippage bounds; min\_out / max\_in parameters & Implemented & Medium \\
\bottomrule
\end{tabularx}
\end{center}

\paragraph{Note on Cairo.} Starknet's execution model differs from EVM but does not eliminate reentrancy-like patterns through external calls and callbacks. Explicit guards and patterns are required.

\subsection{Dependency Risk}

\begin{center}
\begin{tabular}{lll}
\toprule
\textbf{Dependency} & \textbf{Controls} & \textbf{Residual Risk} \\
\midrule
OpenZeppelin Cairo & Pinned versions; monitor security advisories & Upstream vulnerability \\
cairo-fp (fixed-point math) & Extensive precision testing; boundary validation & Math library bugs \\
Starknet protocol & Accept as infrastructure dependency & Protocol-level bugs \\
\bottomrule
\end{tabular}
\end{center}

\section{Upgrade Risk}

All core contracts are upgradeable using OpenZeppelin's Upgradeable component (class hash replacement). This section specifies the upgrade authority, controls, and user protections.

\subsection{Upgrade Impact by Contract}

\begin{center}
\begin{tabular}{lll}
\toprule
\textbf{Contract} & \textbf{Upgrade Can Modify} & \textbf{Impact} \\
\midrule
SY & Deposit/redeem logic; share accounting & Critical \\
PT & Redemption logic; mint authority & Critical \\
YT & Yield distribution; interest accounting & Critical \\
Market & AMM curve; swap logic; fee calculation & Critical \\
Router & Routing logic; could add hidden fees & High \\
Factory / MarketFactory & Implementation templates for new deployments & High \\
\bottomrule
\end{tabular}
\end{center}

\subsection{Upgrade Policy}

\begin{center}
\begin{tabularx}{\textwidth}{l l >{\raggedright\arraybackslash}X l}
\toprule
\textbf{Control} & \textbf{Status} & \textbf{Parameter} & \textbf{Residual} \\
\midrule
Multisig ownership & Planned & Target: 3-of-5 signers & Collusion risk \\
Timelock delay & Planned & Target: 48 hours minimum & Emergency bypass \\
Upgrade events & Implemented & \texttt{Upgraded(class\_hash)} emitted & Requires monitoring \\
Code verification & Required & All class hashes verified on-chain & None \\
Pre-deployment audit & Policy & Audited before mainnet deployment & Audit scope limits \\
\bottomrule
\end{tabularx}
\end{center}

\paragraph{Emergency upgrade policy.}
\begin{itemize}
\item Emergency upgrades (bypassing timelock) are permitted only for active exploits.
\item Emergency actions require documented justification published within 24 hours.
\item Post-emergency review and standard re-deployment within 7 days.
\end{itemize}

\subsection{User Protections During Upgrades}

\begin{center}
\begin{tabularx}{\textwidth}{l >{\raggedright\arraybackslash}X}
\toprule
\textbf{Guarantee} & \textbf{Description} \\
\midrule
Exit window & Timelock provides users N hours to exit before upgrade execution. \\
Always-available operations & Even during pause: \texttt{redeem\_pt\_after\_expiry}, \texttt{burn} (LP withdrawal), \texttt{claim\_interest}. \\
Pausable operations & Swaps and new mints may be paused; existing positions remain accessible. \\
Upgrade transparency & Upgrade proposals include class hash diff and audit link (policy). \\
\bottomrule
\end{tabularx}
\end{center}

\section{Oracle Risk}

The protocol depends on Pragma oracle for yield index data used in YT interest accounting.

\subsection{Oracle Parameters}

\begin{center}
\begin{tabular}{ll}
\toprule
\textbf{Parameter} & \textbf{Value} \\
\midrule
TWAP window & 1 hour (configurable per SY) \\
Staleness threshold & 86400 seconds (24 hours) \\
Data source & Pragma aggregated feed \\
\bottomrule
\end{tabular}
\end{center}

\subsection{Oracle Failure Behavior}

\begin{center}
\begin{tabularx}{\textwidth}{l l >{\raggedright\arraybackslash}X}
\toprule
\textbf{Condition} & \textbf{Protocol Response} & \textbf{Rationale} \\
\midrule
Stale data (age $>$ threshold) & Use watermark index; no new yield distributed & Protects against stale/manipulated data \\
Oracle unavailable & Operations using index continue with last value & Prevents liveness failure \\
Rate below watermark & Watermark unchanged; yield accounting frozen & Prevents negative yield distribution \\
\bottomrule
\end{tabularx}
\end{center}

\paragraph{Operations affected by oracle state.}
\begin{itemize}
\item \texttt{mint\_py}, \texttt{redeem\_py}, \texttt{claim\_interest}: Use current or watermark index.
\item \texttt{redeem\_pt\_after\_expiry}: Does not depend on oracle (1:1 redemption).
\item AMM swaps: Do not use oracle; pricing is internal to pool state.
\end{itemize}

\subsection{Exchange Rate Manipulation}

\paragraph{Threat.} Attacker manipulates oracle to extract value from the protocol.

\paragraph{Attack surfaces.} The oracle index affects only the YT interest accounting system:

\begin{center}
\begin{tabularx}{\textwidth}{l >{\raggedright\arraybackslash}X l}
\toprule
\textbf{Subsystem} & \textbf{Oracle Dependency} & \textbf{Extraction Vector} \\
\midrule
Interest accounting & $I_{global}$ updated from oracle & Inflated index $\Rightarrow$ excess \texttt{claim\_interest} payout \\
SY deposit/redeem & 1:1 ratio (no oracle) & None \\
AMM pricing & Internal pool state only & None (oracle not used) \\
\bottomrule
\end{tabularx}
\end{center}

\paragraph{Attack scenario (index inflation).}
\begin{enumerate}
\item Attacker manipulates oracle to report inflated exchange rate.
\item Attacker calls \texttt{claim\_interest}; inflated $I_{global}$ causes excess SY payout.
\item Oracle returns to normal; attacker has extracted SY from YT contract reserves.
\end{enumerate}

\paragraph{Controls.}
\begin{itemize}
\item TWAP (5-minute window) increases manipulation cost and duration required.
\item Watermark prevents exploitation of rate \emph{drops} (only increases can inflate index).
\item Staleness check rejects data older than 24 hours.
\end{itemize}

\paragraph{Residual risk.} Sustained oracle manipulation over the TWAP window can still inflate the index. Mitigation requires oracle-level security (Pragma's aggregation and data source diversity).

\section{Market Risk}

Market risks affect execution and returns but do not compromise protocol solvency. These are inherent to participation in yield markets.

\subsection{Liquidity Risk}

\begin{center}
\begin{tabularx}{\textwidth}{l >{\raggedright\arraybackslash}X l}
\toprule
\textbf{Scenario} & \textbf{Impact} & \textbf{User Mitigation} \\
\midrule
Thin liquidity & High slippage; inability to exit at fair price & Smaller orders; check price impact \\
LP concentration & Single LP can move market significantly & Monitor LP distribution \\
Extreme proportion & Trades rejected at bounds & Wait for arbitrage; use redemption path \\
\bottomrule
\end{tabularx}
\end{center}

\paragraph{Worst case.} Inability to exit position at fair price before expiry. Protocol solvency is unaffected; this is execution risk only.

\subsection{Impermanent Loss Risk}

IL risk for LPs depends on rate volatility and time to expiry:

\begin{itemize}
\item \textbf{High volatility + long maturity}: Severe IL exposure; rate movements cause significant rebalancing.
\item \textbf{Low volatility or short maturity}: Minimal IL; PT price convergence dominates.
\item \textbf{Near expiry}: IL minimizes as $P_{PT} \to 1$ regardless of rate movements.
\end{itemize}

\paragraph{Key insight.} IL risk decreases as expiry approaches because PT price becomes insensitive to rate changes ($r_{scalar} \to \infty$).

\section{Underlying Asset Risk}

Risks from the yield-bearing assets wrapped by SY tokens.

\begin{center}
\begin{tabularx}{\textwidth}{l >{\raggedright\arraybackslash}X >{\raggedright\arraybackslash}X}
\toprule
\textbf{Risk} & \textbf{Protocol Behavior} & \textbf{User Impact} \\
\midrule
Depeg (value loss) & Watermark prevents negative yield accounting; redemption continues at 1:1 in SY & Principal impairment; users redeem into devalued asset \\
Yield cessation & YT stops accruing; PT converges toward 1 SY & YT holders receive no further yield; PT may reprice \\
Underlying protocol failure & SY may become worthless & Total loss possible \\
Regulatory freeze & Protocol cannot intervene & User bears external legal risk \\
\bottomrule
\end{tabularx}
\end{center}

\paragraph{Key distinction.} The watermark protects yield accounting integrity (no negative yield distribution). It does \emph{not} protect against principal impairment of the underlying asset. If the underlying loses value, all token holders (SY, PT, YT, LP) bear proportional losses.

\section{Governance and Admin Risk}

\subsection{Admin Key Compromise}

\paragraph{Threat.} Attacker gains control of admin keys and executes malicious upgrades.

\paragraph{Impact.} Complete protocol control; ability to drain funds via malicious upgrade.

\begin{center}
\begin{tabularx}{\textwidth}{l l >{\raggedright\arraybackslash}X}
\toprule
\textbf{Control} & \textbf{Status} & \textbf{Protection} \\
\midrule
Hardware wallet storage & Required & Reduces remote key theft \\
Multisig (3-of-5) & Planned & Requires multiple compromises or collusion \\
Timelock (48h) & Planned & Users can exit before malicious upgrade executes \\
Upgrade monitoring & Implemented & On-chain events enable automated alerts \\
\bottomrule
\end{tabularx}
\end{center}

\subsection{Malicious Upgrade Detection}

\begin{center}
\begin{tabularx}{\textwidth}{l >{\raggedright\arraybackslash}X >{\raggedright\arraybackslash}X}
\toprule
\textbf{Attack} & \textbf{Detection Method} & \textbf{User Defense} \\
\midrule
Backdoor insertion & Code diff analysis; class hash verification & Exit during timelock window \\
Hidden fee extraction & Event monitoring; balance tracking & Review upgrade proposals \\
Logic manipulation & Automated contract comparison tools & Subscribe to upgrade alerts \\
\bottomrule
\end{tabularx}
\end{center}

\subsection{Role Escalation}

\begin{center}
\begin{tabular}{ll}
\toprule
\textbf{Risk} & \textbf{Control} \\
\midrule
Admin grants attacker roles & Multisig requirement; role changes emit events \\
Role accumulation & Separation of duties; distinct admin/pauser/operator roles \\
Pauser griefing & Pause is temporary; unpause requires same authority \\
\bottomrule
\end{tabular}
\end{center}

\section{Starknet Infrastructure Risk}

External infrastructure risks outside protocol control.

\begin{center}
\begin{tabularx}{\textwidth}{l >{\raggedright\arraybackslash}X >{\raggedright\arraybackslash}X}
\toprule
\textbf{Risk} & \textbf{Impact} & \textbf{User Response} \\
\midrule
Sequencer downtime & Transactions not processed & Wait for recovery; funds remain in L2 state \\
Sequencer censorship & Specific transactions blocked & Use alternative RPC; await escape mechanisms \\
Proof generation delay & Finality delayed on L1 & Internal accounting unaffected; settlement certainty reduced \\
L1 congestion & High finalization costs & Accept as infrastructure cost \\
Protocol upgrade (Starknet) & Breaking changes possible & Monitor Starknet roadmap; protocol will adapt \\
\bottomrule
\end{tabularx}
\end{center}

\paragraph{Note on censorship.} L1 force-inclusion mechanisms for Starknet are evolving. Until mature escape hatches exist, sequencer censorship remains an infrastructure-level risk that users accept.

\section{User Operational Risk}

User errors and operational security failures.

\begin{center}
\begin{tabularx}{\textwidth}{l >{\raggedright\arraybackslash}X >{\raggedright\arraybackslash}X}
\toprule
\textbf{Risk} & \textbf{Description} & \textbf{Mitigation} \\
\midrule
Wrong address & Funds sent to incorrect recipient & Address validation; UI confirmations; verify full address \\
Address poisoning & Attacker creates similar-looking address & Use address book; verify complete address, not prefix \\
Wrong network/contract & Interaction with fake or test contract & Verify class hash; use official registry \\
Missed expiry & PT held past maturity & Redemption still works; no penalty \\
Slippage misconfiguration & Worse execution than expected & UI defaults; warnings on extreme values \\
Unlimited approvals & Excessive token permissions & Default to exact approvals; revoke after use \\
Approval phishing & Malicious approval requests & Wallet warnings; review approval targets \\
MEV sandwich & Execution between attacker trades & Tight slippage bounds; private mempools if available \\
Private key loss & Permanent fund loss & Hardware wallets; secure backup procedures \\
\bottomrule
\end{tabularx}
\end{center}

\section{Threat Model Summary}

\subsection{Risk Priority Matrix}

Risks are prioritized separately for systemic (protocol-wide) and individual (user-level) impact.

\begin{center}
\begin{tabularx}{\textwidth}{l l l l l}
\toprule
\textbf{Risk} & \textbf{Likelihood} & \textbf{Systemic Impact} & \textbf{User Impact} & \textbf{Priority} \\
\midrule
Admin key compromise & Low & Critical & Critical & Critical \\
Smart contract bug & Low & Critical & Critical & Critical \\
Oracle manipulation & Medium & High & High & High \\
Underlying depeg & Low & High & High & High \\
Liquidity crisis & Medium & Low & Medium & Medium \\
Sequencer downtime & Low & Low & Low & Low \\
User operational error & High & None & High & Medium (user) \\
\bottomrule
\end{tabularx}
\end{center}

\paragraph{Interpretation.}
\begin{itemize}
\item \textbf{Critical priority}: Requires maximum preventive investment and monitoring.
\item \textbf{High priority}: Active controls required; residual risk accepted.
\item \textbf{Medium priority}: Controls in place; ongoing monitoring.
\item \textbf{Low priority}: Accepted as infrastructure or user-level risk.
\end{itemize}

\section{Exclusions: What the Protocol Does NOT Protect Against}

The following risks are explicitly outside the protocol's protection scope:

\begin{enumerate}
\item \textbf{Underlying asset impairment.} Principal loss from underlying asset failure is borne by users.

\item \textbf{Malicious or compromised governance.} Upgrade authority can modify contract logic; mitigated only by timelock and monitoring where implemented.

\item \textbf{Starknet network failures.} Protocol correctness depends on Starknet liveness and validity proofs.

\item \textbf{Oracle compromise.} Fundamental Pragma compromise could cause incorrect yield accounting; mitigated by TWAP and watermark but not eliminated.

\item \textbf{Regulatory actions.} Legal restrictions on assets or users cannot be prevented by the protocol.

\item \textbf{User operational errors.} Incorrect addresses, lost keys, and configuration mistakes are user responsibility.

\item \textbf{MEV and front-running.} Users can mitigate with slippage settings but cannot eliminate; this is inherent to public blockchains.
\end{enumerate}


\chapter{Governance \& Upgrade Model}

This section specifies the governance structure, upgrade capabilities, and admin powers. Most protocol-controlled contracts are upgradeable; PT, YT, Market, RouterStatic, and PyLpOracle are not. Users must evaluate admin trust before depositing.

\paragraph{Authority model.}
All upgradeable contracts use OpenZeppelin's \texttt{OwnableComponent} and \texttt{UpgradeableComponent}. A single authority (the \emph{owner}) holds \texttt{DEFAULT\_ADMIN\_ROLE} and can:
\begin{itemize}
\item Replace any contract's class hash via \texttt{upgrade(new\_class\_hash)}.
\item Grant and revoke roles to other addresses.
\item Execute contract-specific admin functions.
\end{itemize}

\noindent\textbf{Target configuration:} 3-of-5 multisig with 48-hour timelock (not yet deployed).

\paragraph{Upgrade semantics.}
Starknet contracts are upgraded by replacing the \emph{class hash}---the on-chain identifier of the contract's bytecode. When \texttt{upgrade(new\_class\_hash)} executes:
\begin{enumerate}
\item The contract's class hash is atomically replaced.
\item Storage is preserved; the new implementation must be storage-compatible.
\item An \texttt{Upgraded(new\_class\_hash)} event is emitted.
\item All subsequent calls execute the new implementation immediately.
\end{enumerate}
There is no proxy indirection; the contract address remains constant while its code changes.


\section{Contract Classification}

Horizon distinguishes between \emph{internal} contracts (protocol-controlled) and \emph{external dependencies} (third-party infrastructure).

\subsection{Internal Contracts (Protocol-Controlled)}

\begin{center}
\begin{tabularx}{\textwidth}{l l >{raggedright\arraybackslash}X l}
\toprule
\textbf{Contract} & \textbf{Upgradeable} & \textbf{Scope of Upgrade Impact} & \textbf{Risk} \\
\midrule
SY & Yes & Share accounting, deposit/redeem logic & Critical \\
SYWithRewards & Yes & SY accounting plus reward distribution & Critical \\
PT & No & Redemption logic, mint authority validation & Critical \\
YT & No & Yield distribution, interest index accounting & Critical \\
Market & No & AMM curve, swap execution, fee calculation & Critical \\
Factory & Yes & Templates for new SY/PT/YT deployments & High \\
MarketFactory & Yes & Templates for new Market deployments & High \\
Router & Yes & Routing logic, could insert hidden fees & High \\
RouterStatic & No & Read-only quotes and previews & Low \\
PyLpOracle & No & PT/YT/LP TWAP pricing helper & Medium \\
PragmaIndexOracle & Yes & Index source configuration, staleness thresholds & High \\
\bottomrule
\end{tabularx}
\end{center}

\paragraph{Risk interpretation.}
\begin{description}
\item[Critical:] Upgrade can directly drain or freeze user funds.
\item[High:] Upgrade can extract value or disrupt operations but cannot drain existing positions.
\end{description}

\subsection{External Dependencies (Third-Party)}

\begin{center}
\begin{tabular}{l l l}
\toprule
\textbf{Dependency} & \textbf{Trust Requirement} & \textbf{Failure Mode} \\
\midrule
Starknet sequencer & Liveness, censorship resistance & Transactions not processed \\
Pragma oracle & Data integrity, freshness & Incorrect yield accounting \\
Underlying protocols (e.g., Nostra) & Solvency, correct yield accrual & SY value impairment \\
\bottomrule
\end{tabular}
\end{center}

\paragraph{Dependency scope.}
The protocol cannot upgrade or control external dependencies. Users accept these as infrastructure-level trust assumptions. Failures in external dependencies can cause losses even if protocol code is correct.


\section{Admin Powers}

Most admin functions require \texttt{DEFAULT\_ADMIN\_ROLE} unless otherwise noted (e.g., PragmaIndexOracle uses OPERATOR\_ROLE and PAUSER\_ROLE for specific controls). This section enumerates the exact on-chain powers.

\subsection{Powers Available on All Upgradeable Contracts}

Every contract with \texttt{OwnableComponent} and \texttt{UpgradeableComponent} exposes:

\begin{center}
\begin{tabularx}{\textwidth}{l >{raggedright\arraybackslash}X l}
\toprule
\textbf{Function} & \textbf{Effect} & \textbf{Risk} \\
\midrule
\texttt{upgrade(class\_hash)} & Replace contract implementation atomically & Critical \\
\texttt{transfer\_ownership(new\_owner)} & Transfer admin control to new address & Critical \\
\texttt{renounce\_ownership()} & Permanently remove admin control (irreversible) & N/A \\
\bottomrule
\end{tabularx}
\end{center}

\subsection{Contract-Specific Admin Functions}

\paragraph{Factory.}
\begin{center}
\begin{tabularx}{\textwidth}{l >{raggedright\arraybackslash}X l}
\toprule
\textbf{Function} & \textbf{Effect} & \textbf{Risk} \\
\midrule
\texttt{set\_sy\_class\_hash} & Change template for future SY deployments & High \\
\texttt{set\_pt\_class\_hash} & Change template for future PT deployments & Critical \\
\texttt{set\_yt\_class\_hash} & Change template for future YT deployments & Critical \\
\bottomrule
\end{tabularx}
\end{center}

\paragraph{Note.} Class hash changes affect only \emph{new} deployments. Existing SY/PT/YT instances retain their original implementation until individually upgraded.

\paragraph{MarketFactory.}
\begin{center}
\begin{tabularx}{\textwidth}{l >{raggedright\arraybackslash}X l}
\toprule
\textbf{Function} & \textbf{Effect} & \textbf{Risk} \\
\midrule
\texttt{set\_market\_class\_hash} & Change template for future Market deployments & Critical \\
\texttt{set\_treasury} & Update reserve fee recipient & High \\
\texttt{set\_default\_reserve\_fee\_percent} & Configure default reserve fee percentage & Medium \\
\texttt{set\_override\_fee} & Set per-router per-market fee override & Medium \\
\texttt{set\_default\_rate\_impact\_sensitivity} & Configure rate impact fee sensitivity & Medium \\
\texttt{set\_yield\_contract\_factory} & Update PT validation source & Medium \\
\bottomrule
\end{tabularx}
\end{center}

\paragraph{Router.}
\begin{center}
\begin{tabularx}{\textwidth}{l >{raggedright\arraybackslash}X l}
\toprule
\textbf{Function} & \textbf{Effect} & \textbf{Risk} \\
\midrule
\texttt{pause()} & Block all router operations & Medium \\
\texttt{unpause()} & Resume router operations & Low \\
\bottomrule
\end{tabularx}
\end{center}

\paragraph{PragmaIndexOracle.}
\begin{center}
\begin{tabularx}{\textwidth}{l >{raggedright\arraybackslash}X l}
\toprule
\textbf{Function} & \textbf{Effect} & \textbf{Risk} \\
\midrule
\texttt{set\_config} & Update TWAP window and staleness limits & Medium \\
\texttt{pause} / \texttt{unpause} & Freeze/unfreeze oracle reads & Medium \\
\texttt{emergency\_set\_index} & Override stored index (recovery) & High \\
\texttt{initialize\_rbac} & One-time RBAC initialization after upgrade & Medium \\
\bottomrule
\end{tabularx}
\end{center}

\paragraph{Token Contracts (SY, SYWithRewards).}
Token contracts expose only \texttt{upgrade}. Minting/burning authority is enforced by code (YT contract controls PT/YT minting), not by admin-configurable roles.

\paragraph{PT, YT, and Market.}
These contracts are not upgradeable. Their parameters are set at deployment and cannot be modified post-deployment.


\section{Role-Based Access Control}

The protocol uses OpenZeppelin's \texttt{AccessControlComponent} with three roles:

\begin{center}
\begin{tabularx}{\textwidth}{l >{raggedright\arraybackslash}X l}
\toprule
\textbf{Role} & \textbf{Permissions} & \textbf{Holder} \\
\midrule
\texttt{DEFAULT\_ADMIN\_ROLE} & All admin functions; grant/revoke other roles & Owner (multisig target) \\
\texttt{PAUSER\_ROLE} & \texttt{pause()}, \texttt{unpause()} on Router and PragmaIndexOracle & Owner (initially) \\
\texttt{OPERATOR\_ROLE} & PragmaIndexOracle configuration updates & Owner (initially) \\
\bottomrule
\end{tabularx}
\end{center}

\paragraph{Role hierarchy.}
\texttt{DEFAULT\_ADMIN\_ROLE} is the admin of all roles, including itself. Only addresses with \texttt{DEFAULT\_ADMIN\_ROLE} can grant or revoke any role.


\section{Emergency Procedures}

\subsection{Emergency Powers}

The following actions are available during emergencies:

\begin{center}
\begin{tabularx}{\textwidth}{l >{raggedright\arraybackslash}X}
\toprule
\textbf{Action} & \textbf{Effect} \\
\midrule
Pause Router & Blocks new swaps, mints, and complex operations via Router \\
Pause PragmaIndexOracle & Freezes index updates at the stored watermark \\
Emergency upgrade & Replace contract implementation without timelock (if active exploit) \\
\bottomrule
\end{tabularx}
\end{center}

\subsection{Pause Scope}

Pause mechanisms exist on the Router and PragmaIndexOracle. Router pause stops complex protocol interactions, while oracle pause freezes index updates to the stored watermark.

\begin{center}
\begin{tabularx}{\textwidth}{l l >{raggedright\arraybackslash}X}
\toprule
\textbf{Operation} & \textbf{During Pause} & \textbf{Rationale} \\
\midrule
Router swaps & Blocked & Prevent exploitation via complex operations \\
Router minting (PT/YT) & Blocked & Prevent new positions during incident \\
Router liquidity add/remove & Blocked & Prevent position changes \\
\midrule
Direct SY deposit/redeem & Available & Users can exit to underlying \\
Direct PT transfer & Available & Users can move positions \\
Direct YT transfer & Available & Users can move positions \\
Post-expiry PT redemption & Available & Users can redeem matured positions \\
YT interest claim & Available & Users can claim accrued yield \\
Market direct swap & Available & Direct contract calls bypass Router \\
LP burn (withdraw liquidity) & Available & Users can exit LP positions \\
\bottomrule
\end{tabularx}
\end{center}

\begin{center}
\fbox{\parbox{0.9\textwidth}{
\textbf{User Safety Invariant During Emergencies}

Even when the Router is paused, users can always:
\begin{enumerate}
\item Redeem SY for underlying assets (direct call to SY contract).
\item Redeem PT for SY after expiry (direct call to YT contract).
\item Claim accrued YT interest (direct call to YT contract).
\item Withdraw LP positions by burning LP tokens (direct call to Market).
\item Transfer all token types to other addresses.
\end{enumerate}
No emergency action can prevent users from exiting existing positions.
}}
\end{center}

\subsection{Emergency Response Process}

\begin{enumerate}
\item \textbf{Detection.} Exploit identified via monitoring, user report, or external disclosure.
\item \textbf{Triage.} Assess severity and funds at risk. Classify as Critical/High/Medium/Low.
\item \textbf{Containment.} If Critical or High: immediately pause Router.
\item \textbf{Remediation.} Develop and test fix. For Critical exploits, emergency upgrade without timelock.
\item \textbf{Disclosure.} Publish incident report within 24 hours of resolution.
\item \textbf{Post-mortem.} Conduct root cause analysis; implement preventive measures.
\end{enumerate}

\paragraph{Timelock bypass justification.}
Emergency upgrades (bypassing any future timelock) are permitted only when:
\begin{itemize}
\item Active exploitation is occurring or imminent.
\item Funds at risk exceed a material threshold.
\item Waiting for timelock would result in greater user losses.
\end{itemize}


\section{Upgrade Process}

\subsection{Standard Upgrade Steps}

\begin{enumerate}
\item \textbf{Development.} Implement changes; ensure storage compatibility.
\item \textbf{Testing.} Full test suite passes; invariant fuzzing on modified paths.
\item \textbf{Audit.} External audit for significant changes (new features, security-relevant modifications).
\item \textbf{Proposal.} Publish upgrade proposal with:
\begin{itemize}
\item New class hash
\item Code diff or audit report link
\item Rationale and risk assessment
\end{itemize}
\item \textbf{Timelock.} (When implemented) Submit upgrade transaction; wait for timelock period.
\item \textbf{Execution.} Multisig signers approve and execute upgrade.
\item \textbf{Verification.} Confirm new class hash matches proposal; verify contract behavior.
\end{enumerate}

\subsection{Storage Compatibility Requirements}

Upgrades must preserve storage layout:
\begin{itemize}
\item Existing storage variables retain their slots.
\item New variables are appended (not inserted).
\item Variable types are not changed.
\item Mappings and structs maintain compatible layouts.
\end{itemize}

\paragraph{Incompatible upgrades.}
If storage layout must change, a migration pattern is required: deploy new contract, migrate state, update references. This is a breaking change requiring extended user notice.


\section{Governance Attack Surface}

\begin{center}
\begin{tabularx}{\textwidth}{l l l >{raggedright\arraybackslash}X}
\toprule
\textbf{Attack Vector} & \textbf{Likelihood} & \textbf{Impact} & \textbf{Mitigations} \\
\midrule
Admin key compromise & Low & Critical (total loss) & Hardware wallet; multisig (planned); timelock (planned) \\
Malicious insider upgrade & Low & Critical (total loss) & Multisig threshold; code review; audit requirement \\
Social engineering of signers & Medium & Critical & Dedicated signing devices; verification procedures \\
Compromised upgrade proposal & Low & Critical & Class hash verification; multiple reviewer sign-off \\
Role escalation & Very Low & High & RBAC structure; role change events; monitoring \\
Pause abuse (griefing) & Low & Medium (temporary) & PAUSER\_ROLE separation; unpause requires same authority \\
\bottomrule
\end{tabularx}
\end{center}

\paragraph{Attack surface reduction.}
The governance attack surface is proportional to:
\begin{itemize}
\item Number of addresses with admin roles (currently 1; target 5 multisig signers).
\item Frequency of upgrades (minimize by shipping stable code).
\item Complexity of upgrade process (reduce via automation and verification tooling).
\end{itemize}


\section{Trust Assumptions}

Users must evaluate what they trust before using the protocol.

\subsection{Trusted (Cannot Be Verified)}

\begin{itemize}
\item Admin keys are not compromised.
\item Multisig signers do not collude maliciously.
\item Upgrades are deployed with benevolent intent.
\item Emergency actions are taken in good faith.
\item Operational security of key holders is maintained.
\end{itemize}

\subsection{Verifiable On-Chain}

\begin{itemize}
\item All contract code (class hashes can be verified against source).
\item All role grants and revocations (events emitted).
\item All upgrades (events emitted with new class hash).
\item Current pause state.
\item All token balances and positions.
\item AMM state and pool reserves.
\end{itemize}

\subsection{Trustless (Given Current Code)}

\begin{itemize}
\item Mathematical correctness of AMM pricing and yield accounting.
\item Token conservation (no inflation without corresponding deposits).
\item Access control enforcement (only authorized addresses can call admin functions).
\item Expiry mechanics (PT redeems 1:1 after expiry, YT stops accruing).
\end{itemize}

\paragraph{Key insight.}
Upgradeability means ``trustless'' properties hold only for the \emph{current} implementation. A malicious upgrade can violate any property. Users must trust admin or verify each upgrade.


\section{Governance Roadmap}

The following governance features are planned but not yet implemented:

\begin{description}
\item[Multisig (3-of-5).] Require multiple independent signers for all admin actions. Prevents single-key compromise from causing total loss.

\item[Timelock (48 hours).] Delay between upgrade proposal and execution. Provides users time to review changes and exit if concerned.

\item[Upgrade monitoring.] Automated alerts when upgrade transactions are submitted. Enables community review during timelock window.

\item[Verified class hash registry.] On-chain registry of audited, approved class hashes. Upgrades restricted to pre-verified implementations.
\end{description}

\paragraph{Decentralization considerations.}
Full decentralization (governance token, on-chain voting) is not planned for initial deployment. Rationale:
\begin{itemize}
\item Governance tokens introduce new attack vectors (vote buying, governance attacks).
\item Early-stage protocols benefit from rapid iteration capability.
\item Decentralization is meaningful only with sufficient distribution and participation.
\end{itemize}

\noindent The protocol may consider decentralized governance once the codebase is stable and user base is established.


\chapter{Related Work}

This section positions Horizon relative to existing yield and fixed-rate protocols. The comparison focuses on mechanism design, capital efficiency, precision guarantees, and execution environment rather than superficial feature parity.


\section{Pendle Finance (Design Lineage)}

Horizon intentionally adopts Pendle's logit-based AMM and PT/YT conservation model. This is a deliberate design choice: Pendle's mechanics have been battle-tested across multiple market cycles, exhibit well-understood convergence behavior, and create predictable incentives for LPs and arbitrageurs. By preserving design parity, Horizon minimizes economic novelty risk while innovating at the execution and infrastructure layer.

\subsection{Design Parity}

The following elements are intentionally equivalent:

\begin{itemize}
\item \textbf{Logit-based exchange rate:} $\text{rate} = \frac{\ln(p/(1-p))}{r_{scalar}} + r_{anchor}$
\item \textbf{Time-decaying rate scalar:} Sensitivity increases as expiry approaches.
\item \textbf{Conservation invariant:} $\text{PT} + \text{YT} = \text{SY}$ at all times.
\item \textbf{Exponential fee decay:} Fees decrease toward expiry to encourage late-market trading.
\item \textbf{Yield accounting:} Global index tracks cumulative yield; user snapshots determine claimable interest.
\end{itemize}

\paragraph{Rationale.}
Economic equivalence means Horizon inherits Pendle's proven market dynamics: arbitrage bounds implied rates, LPs face predictable IL profiles, and users can reason about positions using existing Pendle mental models. This reduces adoption friction and audit scope.

\subsection{Execution Divergence}

The following elements differ due to execution environment:

\begin{center}
\begin{tabularx}{\textwidth}{l l >{\raggedright\arraybackslash}X}
\toprule
\textbf{Aspect} & \textbf{Pendle (EVM)} & \textbf{Horizon (Starknet)} \\
\midrule
Math library & LogExpMath (Solidity) & cairo-fp 64.64 fixed-point \\
Precision & $\sim$18 decimal digits & $\sim$19 decimal digits (wider intermediates) \\
Overflow model & Manual checks; revert on overflow & Cairo bounds checking; felt252 range \\
Oracle integration & Chainlink + internal rate tracking & Pragma TWAP with configurable staleness \\
Upgrade model & Transparent proxy (EIP-1967) & Class hash replacement (native Starknet) \\
Gas cost per swap & High (L1) / Medium (L2) & Low (validity rollup) \\
\bottomrule
\end{tabularx}
\end{center}

\paragraph{Precision advantage.}
The logit function $\ln(p/(1-p))$ requires high-precision transcendental math. Solidity implementations use lookup tables and approximations; Cairo's 64.64 fixed-point library computes $\ln$ and $\exp$ with lower approximation error, reducing pricing drift at extreme proportions.


\section{Element Finance (Historical Precedent)}

Element Finance pioneered on-chain yield tokenization, demonstrating market demand for PT/YT separation. The protocol operated from 2021--2023 before discontinuing active development.

\subsection{Element's Contribution}

Element validated the core thesis: users want to trade yield separately from principal. The protocol successfully tokenized positions from Yearn vaults, Compound, and Aave, proving that yield tokenization is economically viable.

\subsection{YieldSpace Limitations}

Element used the YieldSpace AMM, a constant-power curve designed for yield trading:

\[
x^{1-t} + y^{1-t} = k
\]

where $t$ represents time to maturity. This curve has known limitations:

\begin{itemize}
\item \textbf{Liquidity fragmentation:} Each expiry requires a separate pool; liquidity cannot be shared.
\item \textbf{Capital inefficiency near maturity:} As $t \to 0$, the curve flattens but does not guarantee $P_{PT} \to 1$.
\item \textbf{Unpredictable convergence:} Final PT price depends on pool composition, not mathematical constraint.
\end{itemize}

\paragraph{Lesson applied.}
Horizon's logit-based curve directly addresses these limitations. The rate scalar $r_{scalar} \to \infty$ as $t \to 0$ mathematically forces $P_{PT} \to 1$, eliminating convergence uncertainty. LPs near expiry face minimal IL because PT price becomes insensitive to trading.


\section{Fixed-Rate Lending Protocols}

Several protocols offer fixed-rate exposure through lending markets rather than yield tokenization. The mechanisms differ fundamentally.

\begin{center}
\begin{tabularx}{\textwidth}{l >{\raggedright\arraybackslash}X >{\raggedright\arraybackslash}X}
\toprule
\textbf{Protocol} & \textbf{Mechanism} & \textbf{Risk Type} \\
\midrule
Notional Finance & Fixed-rate lending via borrow/lend matching & Counterparty risk (borrower default) \\
Yield Protocol & Zero-coupon tokens (fyTokens) via secondary trading & Counterparty risk (undercollateralization) \\
Aave (stable rate) & Variable-to-stable conversion via protocol reserve & Protocol solvency risk \\
\midrule
Horizon & Yield tokenization of existing yield-bearing assets & Underlying asset risk only \\
\bottomrule
\end{tabularx}
\end{center}

\begin{center}
\fbox{\parbox{0.9\textwidth}{
\textbf{Key Economic Distinction}

Horizon does not create yield through leverage or rehypothecation. All returns derive from existing yield-bearing assets (e.g., staked ETH, lending deposits). This eliminates:
\begin{itemize}
\item Borrower default risk (no borrowers exist)
\item Liquidation cascade risk (no collateral to liquidate)
\item Utilization-dependent rate volatility (yield comes from underlying, not borrow demand)
\end{itemize}
PT holders receive exactly 1 SY at expiry regardless of market conditions---the only risk is impairment of the underlying asset itself.
}}
\end{center}

\paragraph{Liquidity risk comparison.}
In lending protocols, liquidity depends on borrow demand and utilization ratios. In yield tokenization, liquidity depends on AMM depth. Both create execution risk, but the risk sources are independent: lending liquidity crises (high utilization) do not affect Horizon markets, and vice versa.


\section{Starknet Execution Environment}

Horizon's design requires specific execution environment properties. This section distinguishes between properties critical to the protocol's mechanics and incidental benefits.

\subsection{Protocol-Critical Properties}

The following Starknet characteristics are necessary for Horizon's design:

\begin{description}
\item[High-precision transcendental math.] The logit function requires $\ln$ and $\exp$ with low approximation error. Cairo's 64.64 fixed-point library provides sufficient precision; EVM alternatives require lookup tables or iterative approximations that accumulate error.

\item[Low-cost arbitrage.] Implied rate accuracy depends on arbitrageurs correcting mispricing. If arbitrage is unprofitable due to gas costs, rates diverge from fair value. Starknet's low transaction costs keep arbitrage profitable even for small deviations.

\item[Affordable oracle updates.] Yield accounting requires frequent index updates. On Ethereum L1, updating yield indices on every relevant transaction would be prohibitively expensive. Starknet's cost structure makes per-transaction oracle reads economically viable.

\item[Native STARK verification.] Starknet's validity proofs ensure transaction correctness without trust assumptions beyond cryptographic soundness. This complements the protocol's trust model.
\end{description}

\subsection{Incidental Benefits}

The following properties improve user experience but are not strictly necessary:

\begin{itemize}
\item \textbf{Fast perceived finality:} L2 confirmation in seconds improves UX for active traders.
\item \textbf{Lower barrier to entry:} Reduced gas costs make small positions economically viable.
\item \textbf{Growing DeFi ecosystem:} Availability of yield-bearing assets (nstSTRK, sSTRK, wstETH) provides initial markets.
\end{itemize}

\paragraph{Alternative environments.}
Other validity rollups (e.g., zkSync, Polygon zkEVM) could theoretically support similar designs. Starknet was chosen for its mature Cairo tooling, native fixed-point libraries, and Pragma oracle integration. The protocol's design is not fundamentally Starknet-exclusive, but porting would require reimplementing the math libraries.


\section{Design References and Prior Work}

The following works inform Horizon's design:

\begin{center}
\begin{tabularx}{\textwidth}{l l >{\raggedright\arraybackslash}X}
\toprule
\textbf{Reference} & \textbf{Year} & \textbf{Design Contribution} \\
\midrule
Pendle Litepaper & 2022 & Logit-based AMM curve; PT/YT conservation model \\
YieldSpace (Yield Protocol) & 2021 & Constant-power curves for yield trading; time-decay mechanics \\
Uniswap v2 Whitepaper & 2020 & Constant product AMM; LP share accounting \\
Curve StableSwap & 2019 & Low-slippage trading for pegged assets; concentrated liquidity \\
Balancer Whitepaper & 2020 & Weighted pools; generalized invariant surfaces \\
\bottomrule
\end{tabularx}
\end{center}

\paragraph{Design lineage.}
Horizon's AMM inherits from Curve's insight that asset-specific invariants outperform general-purpose curves, combined with Pendle's application of this principle to yield markets. The time-decay mechanism originates from YieldSpace but is implemented via the logit transformation rather than constant-power exponents.


\section{Positioning Summary}

Horizon occupies a specific position in the yield protocol landscape:

\begin{center}
\begin{tabularx}{\textwidth}{l >{\raggedright\arraybackslash}X}
\toprule
\textbf{Comparison} & \textbf{Horizon's Position} \\
\midrule
vs Pendle & Equivalent economics; different execution environment (Starknet vs EVM) \\
vs Element & Same tokenization model; superior AMM (logit vs constant-power) \\
vs Fixed-rate lending & Fully asset-backed; no counterparty risk; no leverage \\
vs Generic AMMs & Purpose-built for yield; deterministic expiry convergence \\
\bottomrule
\end{tabularx}
\end{center}

\paragraph{Value proposition.}
Horizon combines Pendle's proven yield tokenization mechanics with Starknet's execution environment to achieve higher arithmetic precision, lower trading friction, and native oracle integration. Compared to fixed-rate lending, Horizon is fully asset-backed and non-leveraged. Compared to earlier yield AMMs, it provides deterministic price convergence and improved capital efficiency near expiry.

The protocol does not claim novel economic mechanisms. Its contribution is bringing battle-tested yield tokenization to a new execution environment with properties---low-cost arbitrage, high-precision math, affordable oracle updates---that enhance the existing design rather than replace it.



\chapter{Conclusion}

\section{Conclusion}

Horizon Protocol introduces a yield tokenization framework on Starknet, enabling users to separate principal and yield exposure for yield-bearing assets.

\paragraph{Key contributions.}
\begin{itemize}
\item \textbf{Starknet-native yield tokenization.} A PT/YT model adapted to Starknet's execution and arithmetic model, providing the first yield tokenization infrastructure on the network.

\item \textbf{Cairo-native numerical design.} High-precision fixed-point arithmetic via the cairo-fp 64.64 library, enabling accurate transcendental function computation for logit-based pricing.

\item \textbf{Native oracle integration.} Pragma TWAP-based exchange rate tracking for yield-bearing assets, with configurable staleness thresholds and watermark protection.

\item \textbf{Logit-based AMM implementation.} A complete implementation of the Pendle-style curve with deterministic price convergence at expiry.
\end{itemize}


\section{Core Guarantees}

The following guarantees define the protocol's economic correctness and distinguish it from lending-based fixed-rate systems.

\begin{center}
\begin{tabular}{l l l}
\toprule
\textbf{Guarantee} & \textbf{Mechanism} & \textbf{Enforcement} \\
\midrule
PT redemption at par & PT $\to$ 1 SY after expiry & Contract logic \\
Token conservation & PT + YT = SY always holds & Arbitrage + contract \\
Deterministic convergence & $r_{scalar} \to \infty$ as $\tau \to 0$ & AMM mathematics \\
Supply equality & PT.supply = YT.supply & Contract assertion \\
Index monotonicity & $I_{global}$ only increases & Watermark logic \\
\bottomrule
\end{tabular}
\end{center}

Together, these guarantees ensure that PT holders receive exactly 1 SY at expiry regardless of market conditions, YT holders receive all accrued yield, and LPs face predictable convergence dynamics.


\section{Design Trade-offs}

The following trade-offs were intentionally chosen to balance correctness, complexity, and deployability for a v1 protocol.

\begin{center}
\begin{tabularx}{\textwidth}{l l >{\raggedright\arraybackslash}X}
\toprule
\textbf{Trade-off} & \textbf{Choice} & \textbf{Rationale} \\
\midrule
Upgradeability vs immutability & Upgradeable & Enables bug fixes; requires admin trust \\
Curve complexity vs simplicity & Full logit curve & Capital efficiency justifies complexity \\
Gas cost vs precision & cairo-fp 64.64 & Financial math requires high precision \\
Single-block vs multi-block execution & Single-transaction & Reduces state complexity; oracle TWAP mitigates manipulation \\
\bottomrule
\end{tabularx}
\end{center}

\paragraph{Note on execution model.}
All protocol operations complete within a single transaction. Multi-block execution patterns (e.g., commit-reveal) were not implemented. Oracle manipulation resistance relies on Pragma's TWAP mechanism rather than protocol-level multi-block designs.


\section{Out-of-Scope Enhancements}

The following enhancements are out of scope for this specification and are not commitments. They are listed to clarify possible future directions.

\begin{center}
\begin{tabular}{l l}
\toprule
\textbf{Enhancement} & \textbf{Purpose} \\
\midrule
Timelock for upgrades & Allow users to exit before upgrade execution \\
Multisig requirement & Reduce single-key compromise risk \\
Additional SY adapters & Support more yield-bearing assets \\
Cross-market routing & Optimize multi-hop trade execution \\
Governance token & Enable decentralized protocol control \\
\bottomrule
\end{tabular}
\end{center}

These items may be addressed in future protocol versions but are explicitly excluded from this specification.


\section{Specification Status}

This document specifies Horizon Protocol v1.0.

\begin{center}
\begin{tabular}{l l l}
\toprule
\textbf{Component} & \textbf{Status} & \textbf{Coverage} \\
\midrule
Token model (SY, PT, YT) & Complete & Mint, redeem, transfer, yield accounting \\
AMM mechanism & Complete & Logit curve, fee decay, liquidity operations \\
Redemption flows & Complete & Pre-expiry and post-expiry paths \\
Oracle integration & Complete & Pragma TWAP, staleness, watermark \\
Governance model & Complete & Upgradeability, RBAC, emergency procedures \\
Risk analysis & Complete & Threat model, attack surfaces, mitigations \\
Parameter specification & Complete & Bounds, defaults, configuration \\
\bottomrule
\end{tabular}
\end{center}

\vspace{1em}

\noindent\textbf{This document constitutes the complete technical specification for Horizon Protocol v1.0. Any behavior not described herein is considered undefined.}



% ============================================================================
% APPENDICES
% ============================================================================

\appendix

\chapter{Mathematical Proofs}

This appendix formalizes the key mathematical properties that underpin Horizon's economic guarantees.


\section{PT Price Convergence}

\textbf{Theorem A.1.}
Assuming a bounded implied rate $r < \infty$, the PT price converges to its redemption value as expiry approaches:
\[
\lim_{\tau \to 0} P_{PT} = 1.
\]

\textbf{Proof.}

The PT price is defined as:
\[
P_{PT} = e^{-r \cdot \tau_y},
\quad
\text{where}
\quad
\tau_y = \frac{\tau}{Y}
\]
and $Y$ is the number of seconds per year (31,536,000).

As $\tau \to 0$, we have $\tau_y \to 0$. Since $r$ is bounded by construction ($0 \le r \le r_{max}$), the exponent satisfies:
\[
\lim_{\tau \to 0} (-r \cdot \tau_y) = 0.
\]

By continuity of the exponential function:
\[
\lim_{\tau \to 0} e^{-r \cdot \tau_y} = e^{0} = 1.
\]

Therefore, $P_{PT} \to 1$ as $\tau \to 0$. \qed


\section{Conservation Invariant}

\textbf{Theorem A.2.}
For all times $t < T$ (before expiry), the following invariant holds:
\[
1 \, \text{PT} + 1 \, \text{YT} \;\longleftrightarrow\; 1 \, \text{SY}.
\]

\textbf{Proof.}

Before expiry, the protocol exposes two mutually inverse state transitions:

\begin{itemize}
\item \textbf{Split:} Depositing $x$ SY mints exactly $x$ PT and $x$ YT.
\item \textbf{Recombine:} Burning $x$ PT and $x$ YT releases exactly $x$ SY.
\end{itemize}

By construction:
\[
\text{PT.supply} = \text{YT.supply}
\]
and all outstanding PT and YT are fully backed by SY held in the YT contract:
\[
\text{SY}_{\text{YT}} = \text{PT.supply} = \text{YT.supply}.
\]

No operation allows PT or YT to be minted or burned independently. Therefore, for all valid states before expiry:
\[
1 \, \text{PT} + 1 \, \text{YT} = 1 \, \text{SY}.
\]

\qed


\section{No-Arbitrage Pricing Condition}

\textbf{Theorem A.3.}
In the absence of fees, slippage, and execution costs, any deviation from
\[
P_{PT} + P_{YT} = 1
\]
admits a risk-free arbitrage opportunity.

\textbf{Assumptions.}
\begin{itemize}
\item Zero swap fees
\item Zero slippage (infinite liquidity)
\item Zero gas/execution costs
\item Instantaneous execution
\end{itemize}

\textbf{Proof.}

\textbf{Case 1:} $P_{PT} + P_{YT} > 1$

\begin{enumerate}
\item Purchase 1 SY at unit price.
\item Split into 1 PT and 1 YT.
\item Sell PT for $P_{PT}$ and YT for $P_{YT}$.
\end{enumerate}

Net profit:
\[
P_{PT} + P_{YT} - 1 > 0.
\]

\textbf{Case 2:} $P_{PT} + P_{YT} < 1$

\begin{enumerate}
\item Purchase 1 PT for $P_{PT}$ and 1 YT for $P_{YT}$.
\item Recombine into 1 SY.
\item Sell SY at unit price.
\end{enumerate}

Net profit:
\[
1 - P_{PT} - P_{YT} > 0.
\]

In both cases, arbitrageurs are incentivized to trade until the equality
\[
P_{PT} + P_{YT} = 1
\]
is restored, up to execution costs.

\qed

\paragraph{Remark.}
In practice, fees and slippage introduce a bounded arbitrage band around the equality. If total round-trip costs are $\epsilon$, then the no-arbitrage condition relaxes to:
\[
|P_{PT} + P_{YT} - 1| \le \epsilon.
\]
Within this band, arbitrage is unprofitable and deviations may persist.



\chapter{Numerical Examples}

The following examples illustrate concrete protocol behavior under simplified assumptions. They are intended for intuition-building and do not model fees, slippage, or adversarial market conditions unless explicitly stated. All calculations use continuous compounding consistent with the protocol's pricing model.


\section{PT Purchase Flow}

\textbf{Scenario.}
A user purchases PT to lock in an implied yield of 10\% APY on 1000 SY, with 180 days to expiry.

\textbf{Assumptions.}
\begin{itemize}
\item No swap fees or slippage
\item Sufficient AMM liquidity (price impact negligible)
\item Continuous compounding, consistent with protocol pricing
\end{itemize}

\textbf{Given:}
\begin{itemize}
\item SY amount: $1000$
\item Implied yield: $r = 0.10$
\item Time to expiry: $\tau_y = \frac{180}{365} \approx 0.493$ years
\end{itemize}

\textbf{PT Price:}
\[
P_{PT} = e^{-r \cdot \tau_y} = e^{-0.10 \times 0.493} \approx 0.9520
\]

\textbf{PT Amount Received:}
\[
\text{PT} = \frac{1000}{0.9520} \approx 1050.4
\]

\textbf{At Maturity:}
\[
\text{SY received} = 1050.4 \times 1 = 1050.4
\]

\textbf{Profit:}
\[
\text{Profit} = 1050.4 - 1000 = 50.4 \text{ SY}
\]

\textbf{Implied APY Check:}
\[
\left(\frac{1050.4}{1000}\right)^{1 / 0.493} - 1 \approx 10.5\%
\]

Note: The discrete APY calculation yields approximately 10.5\% due to the difference between discrete and continuous compounding. Under continuous compounding, the return is exactly 10\% annualized.

\textbf{Interpretation.}
The PT holder sacrifices liquidity and yield upside in exchange for a fixed implied return that is independent of future yield fluctuations. The return is locked at purchase and realized at maturity.


\section{LP Deposit and Withdrawal}

\textbf{Scenario.}
An LP deposits liquidity into a PT/SY market and later withdraws after fees accrue.

\textbf{Assumptions.}
\begin{itemize}
\item PT price remains constant during the period
\item No external trades besides fee-generating swaps
\item PT valued at 1 SY for accounting simplicity (appropriate near expiry)
\end{itemize}

\textbf{Initial Pool State:}
\begin{itemize}
\item Reserves: 1000 SY + 1000 PT
\item Total LP supply: 1000
\end{itemize}

\textbf{Deposit:}
\begin{itemize}
\item LP deposits: 500 SY + 500 PT
\item Deposit ratio: $\frac{500}{1000} = 0.5$
\item LP minted: $0.5 \times 1000 = 500$
\end{itemize}

\textbf{Post-Deposit Pool:}
\begin{itemize}
\item Reserves: 1500 SY + 1500 PT
\item Total LP supply: 1500
\end{itemize}

\textbf{After Fee Accumulation:}
\begin{itemize}
\item Fees earned by pool: 10 SY
\item Pool reserves: 1510 SY + 1500 PT
\end{itemize}

\textbf{Withdrawal (500 LP):}
\[
\text{Ownership share} = \frac{500}{1500} = 0.333
\]

\[
\text{SY out} = 0.333 \times 1510 \approx 503.3
\]
\[
\text{PT out} = 0.333 \times 1500 = 500
\]

\textbf{LP Profit (Fee Income):}
\[
\text{Fee share} = 503.3 - 500 = 3.3 \text{ SY}
\]

The LP receives their pro-rata share of the 10 SY in fees ($\frac{500}{1500} \times 10 = 3.3$ SY).

\textbf{Remark.}
This example isolates fee income by assuming constant PT price. In practice, PT price movements introduce impermanent loss. However, IL risk diminishes as expiry approaches because $P_{PT} \to 1$ regardless of rate movements.


\section{YT Interest Claim}

\textbf{Scenario.}
A user holds YT while the underlying yield index increases.

\textbf{Index Semantics.}
The global index $I_{global}$ represents the cumulative exchange rate of the underlying yield-bearing asset. When the underlying earns yield, $I_{global}$ increases. The user's snapshot $I_{user}$ records the index value at the time of their last claim (or initial acquisition).

\textbf{Given:}
\begin{itemize}
\item YT balance: $B_{YT} = 100$
\item User index (snapshot): $I_{user} = 1.00$
\item Global index (current): $I_{global} = 1.05$
\end{itemize}

\textbf{Interest Calculation:}
\[
\text{Interest} = B_{YT} \times \frac{I_{global} - I_{user}}{I_{user}} = 100 \times \frac{1.05 - 1.00}{1.00} = 5 \text{ SY}
\]

\textbf{Post-Claim State:}
After claiming, the user's index is updated:
\[
I_{user} \leftarrow I_{global} = 1.05
\]

Future claims will only capture yield accrued after this point.

\textbf{Interpretation.}
The YT holder captures the full 5\% increase in the underlying yield index (5 SY on 100 YT). PT holders receive no interim yield---their return is locked at the PT price discount. This separation allows users to speculate on yield rates or hedge yield exposure independently of principal.


\section{Assumptions Summary}

The examples above make the following simplifying assumptions:

\begin{center}
\begin{tabular}{l l}
\toprule
\textbf{Assumption} & \textbf{Real-World Consideration} \\
\midrule
Zero swap fees & Protocol charges fees; reduces effective return \\
Zero slippage & Large trades move price; reduces execution quality \\
Infinite liquidity & AMM has finite depth; large trades have price impact \\
PT price constant & PT price fluctuates with implied rate; creates IL for LPs \\
Instant execution & Block time and sequencer latency introduce timing risk \\
WAD rounding omitted & All amounts are scaled by $10^{18}$; rounding affects dust \\
\bottomrule
\end{tabular}
\end{center}

These assumptions are appropriate for understanding protocol mechanics but should not be used for precise financial projections.



\chapter{Contract Interface Specifications}

This appendix specifies the external interfaces for all core protocol contracts. Each interface includes semantic invariants, parameter specifications, and authority constraints.

\paragraph{Terminology.}
Throughout this appendix, \textbf{PY} denotes a matched pair of PT and YT tokens (1 PY = 1 PT + 1 YT). All amounts are WAD-scaled ($10^{18}$) unless otherwise noted. Timestamps are Unix seconds.


\section{Interface Summary}

\begin{center}
\begin{tabular}{l l l l}
\toprule
\textbf{Interface} & \textbf{Responsibility} & \textbf{Trust Level} & \textbf{Key Invariant} \\
\midrule
SY & Yield abstraction layer & Medium & Fully backed by underlying \\
PT & Principal token claim & High & Supply equals YT supply \\
YT & Yield accounting and distribution & High & Index monotonically increases \\
Market & PT/SY price discovery & Medium & Proportion bounds enforced \\
\bottomrule
\end{tabular}
\end{center}


\section{SY Interface}

\paragraph{Semantic invariants.}
\begin{itemize}
\item SY is always fully backed by underlying assets held by the contract.
\item \texttt{exchange\_rate()} is monotonically non-decreasing (yield only accrues, never decreases).
\item Total SY supply equals the underlying value divided by exchange rate.
\end{itemize}

\begin{lstlisting}[language=C,basicstyle=\ttfamily\small,frame=single,breaklines=true,columns=fullflexible]
#[starknet::interface]
pub trait ISY<TContractState> {
    // ERC20 standard functions (name, symbol, decimals,
    // total_supply, balance_of, transfer, approve, transfer_from)

    fn deposit(
        ref self: TContractState,
        receiver: ContractAddress,
        token_in: ContractAddress,
        amount_token_to_deposit: u256,
        min_shares_out: u256
    ) -> u256;

    fn redeem(
        ref self: TContractState,
        receiver: ContractAddress,
        amount_shares_to_redeem: u256,
        token_out: ContractAddress,
        min_token_out: u256,
        burn_from_internal_balance: bool
    ) -> u256;

    fn exchange_rate(self: @TContractState) -> u256;
    fn underlying(self: @TContractState) -> ContractAddress;
    fn yt(self: @TContractState) -> ContractAddress;
}
\end{lstlisting}

\paragraph{Parameter specification.}
\begin{center}
\begin{tabular}{l l l}
\toprule
\textbf{Parameter} & \textbf{Type} & \textbf{Description} \\
\midrule
\texttt{amount\_token\_to\_deposit} & u256 & Underlying amount (WAD-scaled) \\
\texttt{min\_shares\_out} & u256 & Minimum SY to receive; reverts if not met \\
\texttt{amount\_shares\_to\_redeem} & u256 & SY amount to burn (WAD-scaled) \\
\texttt{token\_in / token\_out} & address & Underlying token address (multi-asset support) \\
\texttt{exchange\_rate} return & u256 & Current rate (WAD-scaled); $\geq 1 \times 10^{18}$ \\
\bottomrule
\end{tabular}
\end{center}

\paragraph{Notes.}
\begin{itemize}
\item \texttt{token\_in} / \texttt{token\_out} enable SY implementations that support multiple underlying assets (e.g., wrapped vs unwrapped).
\item \texttt{burn\_from\_internal\_balance} allows gas-efficient redemptions when SY is already held by the contract.
\item \texttt{deposit} reverts if \texttt{min\_shares\_out} is not met (slippage protection).
\end{itemize}


\section{PT Interface}

\paragraph{Semantic invariants.}
\begin{itemize}
\item PT is non-inflationary: the only mint/burn authority is the YT contract.
\item PT supply always equals YT supply.
\item After expiry, each PT is redeemable for exactly 1 SY.
\end{itemize}

\begin{lstlisting}[language=C,basicstyle=\ttfamily\small,frame=single,breaklines=true,columns=fullflexible]
#[starknet::interface]
pub trait IPT<TContractState> {
    // ERC20 standard functions (inherited)

    fn expiry(self: @TContractState) -> u64;
    fn is_expired(self: @TContractState) -> bool;
    fn sy(self: @TContractState) -> ContractAddress;
    fn yt(self: @TContractState) -> ContractAddress;

    // Only callable by YT contract
    fn mint(ref self: TContractState, to: ContractAddress, amount: u256);
    fn burn(ref self: TContractState, from: ContractAddress, amount: u256);
}
\end{lstlisting}

\paragraph{Parameter specification.}
\begin{center}
\begin{tabular}{l l l}
\toprule
\textbf{Parameter} & \textbf{Type} & \textbf{Description} \\
\midrule
\texttt{expiry} return & u64 & Unix timestamp (seconds) of maturity \\
\texttt{is\_expired} return & bool & True if current time $\geq$ expiry \\
\texttt{amount} (mint/burn) & u256 & PT amount (WAD-scaled) \\
\bottomrule
\end{tabular}
\end{center}

\paragraph{Access control.}
\texttt{mint} and \texttt{burn} are restricted to the YT contract address. Calls from any other address revert with ``unauthorized''.


\section{YT Interface}

\paragraph{Semantic invariants.}
\begin{itemize}
\item YT supply always equals PT supply.
\item The global index $I_{global}$ only increases (watermark protection).
\item Interest is distributed proportionally based on YT balance and index delta.
\item After expiry, YT stops accruing yield and becomes non-transferable.
\end{itemize}

\paragraph{Interest accounting model.}
\begin{itemize}
\item Interest accrues via a global monotonic index updated from the oracle.
\item Each user has a snapshot index $I_{user}$ recording their last claim.
\item Claimable interest: $B_{YT} \times (I_{global} - I_{user}) / I_{user}$.
\item Index updates occur on all state-changing operations.
\end{itemize}

\begin{lstlisting}[language=C,basicstyle=\ttfamily\small,frame=single,breaklines=true,columns=fullflexible]
#[starknet::interface]
pub trait IYT<TContractState> {
    // ERC20 standard functions (inherited)

    // Minting (splits SY into PT + YT)
    fn mint_py(
        ref self: TContractState,
        receiver: ContractAddress,
        amount_sy_to_mint: u256
    ) -> (u256, u256);  // returns (pt_amount, yt_amount)

    // Redemption (recombines PT + YT into SY)
    fn redeem_py(
        ref self: TContractState,
        receiver: ContractAddress,
        amount_py_to_redeem: u256
    ) -> u256;  // returns sy_amount

    fn redeem_py_and_underlying(
        ref self: TContractState,
        receiver: ContractAddress,
        amount_py_to_redeem: u256,
        token_out: ContractAddress
    ) -> (u256, u256);  // returns (sy_redeemed, underlying_out)

    // Interest
    fn claim_interest(ref self: TContractState, user: ContractAddress) -> u256;
    fn get_user_interest(self: @TContractState, user: ContractAddress) -> u256;

    // State
    fn expiry(self: @TContractState) -> u64;
    fn is_expired(self: @TContractState) -> bool;
    fn py_index_stored(self: @TContractState) -> u256;
    fn sy(self: @TContractState) -> ContractAddress;
    fn pt(self: @TContractState) -> ContractAddress;
}
\end{lstlisting}

\paragraph{Parameter specification.}
\begin{center}
\begin{tabular}{l l l}
\toprule
\textbf{Parameter} & \textbf{Type} & \textbf{Description} \\
\midrule
\texttt{amount\_sy\_to\_mint} & u256 & SY to split into PT+YT (WAD-scaled) \\
\texttt{amount\_py\_to\_redeem} & u256 & PT+YT pairs to recombine (WAD-scaled) \\
\texttt{py\_index\_stored} return & u256 & Current global index (WAD-scaled) \\
\texttt{claim\_interest} return & u256 & SY amount transferred to user \\
\bottomrule
\end{tabular}
\end{center}

\paragraph{Notes.}
\begin{itemize}
\item \texttt{mint\_py} requires prior SY approval; transfers SY from caller and mints equal PT and YT to receiver.
\item \texttt{redeem\_py} requires caller to hold equal amounts of PT and YT; burns both and returns SY.
\item \texttt{claim\_interest} updates $I_{user}$ to $I_{global}$ after transfer; subsequent calls return 0 until new yield accrues.
\item After expiry, \texttt{redeem\_pt\_after\_expiry} (not shown) allows PT-only redemption at 1:1.
\end{itemize}


\section{Market Interface}

\paragraph{Semantic invariants.}
\begin{itemize}
\item Pool proportions remain within configured bounds; trades violating bounds revert.
\item LP tokens represent pro-rata ownership of pool reserves.
\item Swap prices follow the logit-based AMM curve (see Section 4).
\item Fees accrue to the pool and are distributed to LPs on withdrawal.
\end{itemize}

\begin{lstlisting}[language=C,basicstyle=\ttfamily\small,frame=single,breaklines=true,columns=fullflexible]
#[starknet::interface]
pub trait IMarket<TContractState> {
    // ERC20 for LP tokens (inherited)

    // Swaps
    fn swap_exact_pt_for_sy(
        ref self: TContractState,
        receiver: ContractAddress,
        exact_pt_in: u256,
        min_sy_out: u256
    ) -> u256;  // returns sy_out

    fn swap_exact_sy_for_pt(
        ref self: TContractState,
        receiver: ContractAddress,
        exact_sy_in: u256,
        min_pt_out: u256
    ) -> u256;  // returns pt_out

    // Liquidity
    fn add_liquidity(
        ref self: TContractState,
        receiver: ContractAddress,
        sy_desired: u256,
        pt_desired: u256,
        min_lp_out: u256
    ) -> (u256, u256, u256);  // returns (sy_used, pt_used, lp_minted)

    fn remove_liquidity(
        ref self: TContractState,
        receiver: ContractAddress,
        lp_to_burn: u256,
        min_sy_out: u256,
        min_pt_out: u256
    ) -> (u256, u256);  // returns (sy_out, pt_out)

    // State
    fn get_reserves(self: @TContractState) -> (u256, u256);
    fn get_ln_implied_rate(self: @TContractState) -> u256;
    fn expiry(self: @TContractState) -> u64;
    fn pt(self: @TContractState) -> ContractAddress;
    fn sy(self: @TContractState) -> ContractAddress;
}
\end{lstlisting}

\paragraph{Parameter specification.}
\begin{center}
\begin{tabular}{l l l}
\toprule
\textbf{Parameter} & \textbf{Type} & \textbf{Description} \\
\midrule
\texttt{exact\_pt\_in / exact\_sy\_in} & u256 & Exact input amount (WAD-scaled) \\
\texttt{min\_sy\_out / min\_pt\_out} & u256 & Minimum output; reverts if not met \\
\texttt{sy\_desired / pt\_desired} & u256 & Maximum amounts to deposit \\
\texttt{min\_lp\_out} & u256 & Minimum LP tokens; reverts if not met \\
\texttt{get\_reserves} return & (u256, u256) & Current (SY, PT) reserves \\
\texttt{get\_ln\_implied\_rate} return & u256 & $\ln(1+r)$ scaled by WAD \\
\bottomrule
\end{tabular}
\end{center}

\paragraph{Pricing reference.}
All swap prices follow the logit-based AMM defined in Section 4. The exchange rate is computed as:
\[
\text{rate} = \frac{\ln(p/(1-p))}{r_{scalar}} + r_{anchor}
\]
where $p$ is the PT proportion and $r_{scalar}$ decays toward expiry.

\paragraph{Slippage protection.}
All swap and liquidity functions include \texttt{min\_*\_out} parameters. If the computed output is less than the minimum, the transaction reverts. This protects users against price movement between transaction submission and execution.


\section{Revert Conditions Summary}

\begin{center}
\begin{tabular}{l l}
\toprule
\textbf{Condition} & \textbf{Affected Functions} \\
\midrule
Output below minimum & All swaps, deposits, withdrawals \\
Insufficient balance & All transfers, burns, redemptions \\
Unauthorized caller & PT mint/burn (non-YT caller) \\
Proportion bounds violated & Market swaps \\
Zero amount & All state-changing functions \\
Expired (where applicable) & \texttt{mint\_py}, \texttt{redeem\_py} \\
Not expired (where applicable) & \texttt{redeem\_pt\_after\_expiry} \\
\bottomrule
\end{tabular}
\end{center}



\chapter{Glossary}
\addcontentsline{toc}{chapter}{Glossary}

\section*{Glossary}
\label{sec:glossary}

This glossary defines key terms used throughout this specification. Terms are ordered alphabetically.

\begin{center}
\begin{tabularx}{\textwidth}{l >{\raggedright\arraybackslash}X}
\toprule
\textbf{Term} & \textbf{Definition} \\
\midrule

\textbf{AMM} &
Automated Market Maker; a smart contract that provides liquidity and determines prices algorithmically based on pool reserves and a pricing function. \\

\textbf{cairo-fp} &
Cairo fixed-point arithmetic library using 64.64 format (64 integer bits, 64 fractional bits), enabling high-precision transcendental function computation. \\

\textbf{Expiry ($T$)} &
Unix timestamp (seconds) at which PT becomes redeemable for exactly 1 SY and YT stops accruing yield. \\

\textbf{Implied Rate ($r$)} &
The annualized yield rate inferred from the current PT/SY exchange rate, assuming PT converges to 1 SY at expiry. Computed as $r = -\ln(P_{PT}) / \tau_y$. \\

\textbf{Logit} &
Transformation $\ln(p/(1-p))$ mapping pool proportions in $(0,1)$ to the real line $(-\infty, +\infty)$, used to derive the AMM exchange rate. \\

\textbf{LP Token} &
Token representing proportional ownership of a PT/SY liquidity pool; entitles holder to pro-rata share of reserves and accumulated fees. \\

\textbf{Market} &
Smart contract implementing the PT/SY AMM, responsible for swaps, liquidity provision, and price discovery. \\

\textbf{Minimum Liquidity} &
Small amount of LP tokens (typically $10^3$) permanently locked on first deposit to prevent share inflation attacks. \\

\textbf{PT (Principal Token)} &
Token representing a claim on principal that redeems for exactly 1 unit of SY at or after expiry. Trades at a discount before expiry reflecting the implied yield. \\

\textbf{PY} &
A matched pair of PT and YT tokens (1 PY = 1 PT + 1 YT), which together are economically equivalent to 1 SY before expiry. \\

\textbf{PY Index ($I_{global}$)} &
Monotonically non-decreasing cumulative index tracking yield accrual for YT holders. Updated from oracle on state-changing operations. \\

\textbf{Rate Anchor ($r_{anchor}$)} &
Dynamic offset parameter that ensures implied yield continuity across successive trades; updated after each swap to preserve rate consistency. \\

\textbf{Rate Scalar ($r_{scalar}$)} &
Time-dependent scaling parameter that controls AMM price sensitivity to changes in pool composition. Increases toward infinity as expiry approaches, forcing $P_{PT} \to 1$. \\

\textbf{RBAC} &
Role-Based Access Control system governing privileged contract operations (upgrades, pausing, parameter changes). \\

\textbf{SY (Standardized Yield)} &
Wrapper token that standardizes heterogeneous yield-bearing assets into a uniform accounting and redemption interface. Implements ERC-20 with deposit/redeem functions. \\

\textbf{Time to Expiry ($\tau$)} &
Seconds remaining until expiry: $\tau = T - t_{now}$. Often normalized to years: $\tau_y = \tau / Y$ where $Y = 31{,}536{,}000$. \\

\textbf{TWAP} &
Time-Weighted Average Price; oracle mechanism that averages prices over a time window to mitigate short-term manipulation. \\

\textbf{Underlying Asset} &
The original yield-bearing asset (e.g., nstSTRK, wstETH) that is wrapped by SY. Distinct from SY itself. \\

\textbf{WAD} &
Fixed-point scaling factor equal to $10^{18}$, used for all external-facing token amounts, rates, and prices. \\

\textbf{Watermark} &
Accounting mechanism that prevents the yield index from decreasing, protecting YT holders' accrued yield from negative oracle movements. \\

\textbf{YT (Yield Token)} &
Token representing the exclusive right to all yield generated by the underlying asset from mint time until expiry. Yield is claimable via \texttt{claim\_interest}. \\

\bottomrule
\end{tabularx}
\end{center}


\section*{Symbol Reference}

\begin{center}
\begin{tabular}{l l l}
\toprule
\textbf{Symbol} & \textbf{Name} & \textbf{Units} \\
\midrule
$P_{PT}$ & PT price in SY terms & Dimensionless (0 to 1) \\
$P_{YT}$ & YT price in SY terms & Dimensionless (0 to 1) \\
$r$ & Implied annual rate & Per-year (dimensionless) \\
$\tau$ & Time to expiry & Seconds \\
$\tau_y$ & Time to expiry (annualized) & Years \\
$T$ & Expiry timestamp & Unix seconds \\
$Y$ & Seconds per year & 31,536,000 \\
$I_{global}$ & Global yield index & WAD-scaled ratio \\
$I_{user}$ & User yield index snapshot & WAD-scaled ratio \\
$r_{scalar}$ & Rate scalar parameter & WAD-scaled \\
$r_{anchor}$ & Rate anchor parameter & WAD-scaled \\
$p$ & PT proportion in pool & Dimensionless (0 to 1) \\
\bottomrule
\end{tabular}
\end{center}

\vspace{2em}

\begin{center}
\rule{0.5\textwidth}{0.4pt}
\end{center}

\vspace{1em}

\begin{center}
\textit{End of Specification}
\end{center}

\vspace{1em}

\noindent\textbf{Document Version:} v2.0\\
\textbf{Last Updated:} December 2025\\
\textbf{License:} BUSL-1.1 (converts to GPL-3.0 on 2028-12-19)

\vspace{1em}

\noindent\textit{This document is a technical specification and does not constitute investment, legal, or financial advice. Users should conduct independent analysis before interacting with the protocol.}


% ============================================================================
% DOCUMENT END
% ============================================================================

\end{document}
